\documentclass[12pt,a4paper]{report}

\usepackage[utf8]{inputenc}
\usepackage[T1]{fontenc}
\usepackage{lmodern}
\usepackage{microtype}
\usepackage{geometry}
\geometry{a4paper, margin=2.5cm, headheight=14pt}
\usepackage{setspace}
\onehalfspacing
\usepackage{amsmath,amssymb,amsthm,mathtools}
\usepackage{booktabs}
\usepackage{array}
\usepackage{multirow}
\usepackage{enumitem}
\usepackage{caption}
\usepackage{subcaption}
\usepackage{algorithm}
\usepackage{algpseudocode}
\usepackage{tikz}
\usetikzlibrary{arrows.meta,positioning,fit,calc,shapes.geometric,backgrounds}
\usepackage{pgfplots}
\pgfplotsset{compat=1.18}
\usepackage[numbers,sort&compress]{natbib}
\usepackage[hidelinks]{hyperref}
\makeatletter
\renewcommand{\theALG@line}{\arabic{algorithm}.\arabic{ALG@line}}
\AtBeginDocument{\renewcommand{\theHALG@line}{\arabic{algorithm}.\arabic{ALG@line}}}
\makeatother
\usepackage[nameinlink,noabbrev]{cleveref}

\newtheorem{definition}{Definition}[chapter]
\newtheorem{assumption}{Assumption}[chapter]
\newtheorem{proposition}{Proposition}[chapter]

\DeclareMathOperator*{\argmax}{arg\,max}
\DeclareMathOperator{\round}{round}
\DeclareMathOperator{\median}{median}
\DeclareMathOperator{\softmax}{softmax}
\DeclareMathOperator{\diag}{diag}

\setlist[itemize]{topsep=0.35em,itemsep=0.2em}
\setlist[enumerate]{topsep=0.35em,itemsep=0.2em}
\captionsetup{font=small,labelfont=bf,skip=6pt}
\sloppy
\hbadness=10000
\emergencystretch=2em

\title{Research on an Automated Box-Counting Method Using a Deep Instance Segmentation Model}
\author{}
\date{}

\begin{document}

\begin{titlepage}
\hypersetup{pageanchor=false}
\centering
\vspace*{1.2cm}
{\large Computer Science Thesis}\\[0.35cm]
{\large System Design and Optimization}\\[1.6cm]
{\LARGE\bfseries Research on an Automated Box-Counting Method Using a Deep Instance Segmentation Model\par}
\vspace{1.1cm}
{\large An update of the YOLO-vf-seg instance segmenter for side-view counting and color typing of beer crates on a forklift\par}
\vspace{1.4cm}
\begin{tikzpicture}[scale=0.92, every node/.style={font=\small}]
  \draw[thick] (0,0) -- (11.2,0);
  \draw[fill=black!8] (0.3,0.25) rectangle (2.5,1.35);
  \draw[fill=black!18] (0.7,1.35) rectangle (1.9,2.15);
  \draw[line width=1pt] (2.5,0.38) -- (4.15,0.38);
  \draw[line width=1pt] (2.5,0.58) -- (4.15,0.58);
  \draw[fill=brown!45] (4.15,0.22) rectangle (8.55,0.72);
  \foreach \c in {0,1,2} {
    \foreach \r in {0,1,2} {
      \draw[fill=orange!35, draw=black!70] ({4.35+\c*1.35},{0.82+\r*0.78}) rectangle ++(1.22,0.68);
    }
  }
  \draw[fill=green!35, draw=black!70] (4.35,3.16) rectangle ++(1.22,0.68);
  \draw[fill=red!35, draw=black!70] (5.70,3.16) rectangle ++(1.22,0.68);
  \draw[fill=blue!30, draw=black!70] (7.05,3.16) rectangle ++(1.22,0.68);
  \draw[fill=black!75] (9.55,2.15) rectangle (10.55,2.62);
  \draw[-{Stealth}, thick] (9.55,2.38) -- (8.7,2.15);
  \node at (1.4,0.8) {\footnotesize forklift};
  \node at (6.35,0.47) {\footnotesize pallet};
  \node[align=center] at (10.05,3.05) {\footnotesize side\\[-2pt]\footnotesize camera};
\end{tikzpicture}
\vfill
{\large November 2026}
\hypersetup{pageanchor=true}
\end{titlepage}

\pagenumbering{roman}
\begin{abstract}
\noindent
This thesis studies the last visual check of a brewery packaging line: beer crates stacked on a pallet, carried by a forklift into the warehouse. A camera looking at the side of the load must name each crate by body color and must count how many crates pass. The stack is regular, the colors are few, and most of the geometry is known in advance. The image is not. Forks shadow the bottom row, the top row is often cut by the frame, adjacent crates of the same color share a weak seam, and the aisle mixes daylight with lamps whose color changes through the day.

The starting point is YOLO-vf-seg, an instance segmenter in the one-stage YOLO family whose neck is biased toward the horizontal lips of a side-view stack. Two updates are developed on top of that segmenter. The first, the illumination-invariant chromatic prototype head (ICPH), does not ask a global white-balance algorithm to repaint the frame. It builds, inside each predicted mask, a feature that keeps reflectance chromaticity, opponent color, and an estimate of the local illuminant, and it scores that feature against a learned prototype of each crate color. The second, perspective lattice completion (PLC), rectifies the front face, snaps mask centers to the known crate pitch, fills a single-cell gap only when orthogonal neighbors support it, and multiplies the occupied front cells by a pallet-depth reading.

Both updates are estimated by a Monte Carlo observation model, not by a recording from a brewery aisle. Under that model, ICPH holds color accuracy better than raw prototype matching once the illuminant becomes strong, and it avoids the collapse that gray-world and shades-of-gray corrections produce on a nearly uniform plastic face. PLC, applied to one frame, is the most accurate of the six counters that were compared. Fusing every frame of a passage by a simple union is not: false cells accumulate and the gap rule then invents crates. The thesis therefore specifies the system around a single sharp frame, states the geometry and the timing budget, and treats the numerical chapter as an estimate of the two modules under an explicit generative model.
\end{abstract}

\tableofcontents
\listoffigures
\listoftables

\chapter*{Principal notation}
\addcontentsline{toc}{chapter}{Principal notation}
\markboth{Principal notation}{Principal notation}

The notation below is used from Chapter~2 onward. Earlier sections introduce each symbol again at the point where it is needed, so that the survey chapter can be read on its own.

\begin{center}
\begin{tabular}{@{}ll@{}}
\toprule
Symbol & Meaning \\
\midrule
$I_c(x)$ & Intensity of channel $c\in\{r,g,b\}$ at pixel $x$ \\
$\mathbf{L}$ & Diagonal illuminant, three positive gains \\
$\mathbf{R}(x)$ & Surface reflectance at $x$ \\
$\Omega$, $M$ & A crate face, and the predicted mask of one instance \\
$p$ & Order of the Minkowski illuminant estimate \\
$\hat{\ell}$ & Estimated illuminant on a mask \\
$\chi$ & Chromaticity of the corrected face \\
$o$ & Two-dimensional opponent color \\
$\eta$ & Chromaticity of the illuminant estimate \\
$\mathbf{f}$ & Eight-dimensional ICPH feature \\
$W,b$ & Linear prototype classifier \\
$w,h$ & Nominal width and height of the front face of one crate \\
$d$, $\hat{d}$ & True pallet depth in crates, and the depth reading \\
$(r,c)$ & Integer lattice cell, row then column \\
$\mathcal{F}$ & Set of front cells declared occupied \\
$\hat{N}$ & Estimated crate count, $\lvert\mathcal{F}\rvert\cdot\hat{d}$ \\
$H$ & Homography from the image to the front plane \\
\bottomrule
\end{tabular}
\end{center}

\vspace{0.6em}
\noindent
Nominal numbers used throughout the design and the estimate are a face of $w=0.40\,\mathrm{m}$ by $h=0.28\,\mathrm{m}$, a crate depth of $0.30\,\mathrm{m}$, six body colors (amber, green, red, blue, yellow, white), and a passage of five frames when a temporal variant is tested. These are design values for returnable plastic crates of one family. They are not claimed as a measurement of any particular brewery's crate.

\cleardoublepage
\pagenumbering{arabic}

\chapter{Related Work and the Problem This Thesis Addresses}
\label{ch:related}

The counting problem in this thesis sits at the junction of three older lines of work. The first is instance segmentation: the assignment of a separate mask to every object, rather than a single label per pixel. The second is visual counting, which began with crowds and cells and has only recently been asked to count goods that hide one another. The third is computational color constancy, the attempt to separate the color of a surface from the color of the light that falls on it. A brewery aisle forces all three to be solved together, on a moving forklift, from a camera that sees only one face of the stack. This chapter sets out the principles, then the systems that have already been built for warehouses and cartons, and then the precise gap that the rest of the thesis fills.

\section{Principles of segmentation, counting, and color}
\label{sec:principles}

\subsection{What the camera records}
\label{sec:formation}

A digital camera does not record the paint on a crate. It records a number in each of three channels, and that number mixes the spectrum of the lamp, the reflectance of the plastic, the spectral sensitivity of the sensor, and the exposure chosen by the shutter. Under the standard simplification used throughout color constancy, a Lambertian surface under a single illuminant collapses to a diagonal model~\cite{gijsenij2011survey}. If $\mathbf{R}(x)\in\mathbb{R}^{3}$ is the reflectance of the surface at pixel $x$ and $\mathbf{L}\in\mathbb{R}^{3}$ is a positive illuminant vector, the observed intensity is the element-wise product
\begin{equation}
\mathbf{I}(x) = \mathbf{L} \odot \mathbf{R}(x).
\label{eq:diagonal}
\end{equation}
The same crate, moved from a north-facing door to a sodium lamp deeper in the aisle, changes $\mathbf{L}$ and therefore changes $\mathbf{I}$, while $\mathbf{R}$ stays fixed. Any classifier that memorizes $\mathbf{I}$ is memorizing the lamp as well as the plastic.

Three further effects sit on top of~\eqref{eq:diagonal} in the scene of this thesis, and none of them is a small correction. Shading from the forks and from the crate above multiplies $\mathbf{I}$ by a scalar that varies smoothly across the face. A neighboring crate of another color leaks into the mask wherever the predicted boundary is a few pixels too wide. Motion of the forklift, if the shutter is open for more than a few milliseconds, smears the horizontal lip of each crate along the direction of travel. The diagonal model remains the right first equation because the two updates in Chapter~\ref{ch:methods} attack illuminant and geometry separately, but the estimate in Chapter~\ref{ch:system} puts shading, edge mixing, and a detection loss that grows with blur back into the observation model so that the diagonal story is not asked to carry the whole scene.

The geometry of the same image is a projective map~\cite{hartley2004mvg}. A pinhole camera sends a point $\mathbf{X}$ on the front face of the stack to a pixel
\begin{equation}
s \begin{bmatrix} u \\ v \\ 1 \end{bmatrix}
= K \,[\,R \mid \mathbf{t}\,]\,
\begin{bmatrix} \mathbf{X} \\ 1 \end{bmatrix},
\label{eq:pinhole}
\end{equation}
with $K$ the intrinsics, $R$ and $\mathbf{t}$ the pose of the camera relative to the face, and $s$ a scale. When the front face is approximately a plane, the map from that plane to the image is a homography. Calibration of $K$ by a planar target~\cite{zhang2000calibration} and estimation of the homography are therefore sufficient to put detections into metres on the crate face. That metric plane is what the lattice counter of Chapter~\ref{ch:methods} assumes. It is not a reconstruction of the whole pallet in three dimensions.

\subsection{Detection, semantic segmentation, and instance segmentation}
\label{sec:tasks}

Object detection returns a box and a class. The modern two-stage line begins with region proposals and a convolutional classifier~\cite{girshick2014rcnn}, speeds the classifier by sharing the feature map~\cite{girshick2015fastrcnn}, and then learns the proposals themselves~\cite{ren2015faster}. A cascade of detectors trained at increasing intersection-over-union tightens the boxes that survive~\cite{cai2018cascade}. One-stage detectors skip the proposal stage. SSD places default boxes on several feature maps~\cite{liu2016ssd}. RetinaNet shows that the extreme imbalance between background anchors and object anchors can be handled by a focal loss that down-weights the easy negatives~\cite{lin2017focal}. FCOS removes anchors and regresses a box from each foreground pixel~\cite{tian2019fcos}. CornerNet detects a pair of corners and groups them~\cite{law2018cornernet}. CenterNet detects the center point and reads the box size from that point~\cite{zhou2019centernet}. YOLO, from its first version, frames detection as one regression on a grid~\cite{redmon2016yolo}.

Semantic segmentation assigns a class to every pixel and does not say which crate is which. Fully convolutional networks showed that a classifier can be turned into a dense predictor~\cite{long2015fcn}. U-Net added a skip-connected decoder that restores spatial detail~\cite{ronneberger2015unet}. DeepLabv3+ combined atrous convolution in the encoder with a light decoder~\cite{chen2018deeplab}. None of these outputs can be counted by connected components when two crates of the same color touch, because the two crates are one component.

Instance segmentation returns a separate mask per object. Mask R-CNN adds a small fully convolutional mask head on each region of interest of Faster R-CNN, and aligns the pooled features with the box by bilinear sampling~\cite{he2017maskrcnn}. The mask head is class-specific in the original design: the network predicts a mask for every class and the box class selects which mask is kept. That separation matters for the brewery. Color type and spatial extent are different questions. A mask can be right while the color name is wrong, and the counting stage should be able to use the mask either way.

The optimization and generalization background used throughout is the standard one for deep networks~\cite{goodfellow2016dl,lecun2015nature}. Training these models depends on large labeled sets. Pascal VOC established the detection challenge that the first deep detectors were measured on~\cite{everingham2010pascal}. ImageNet supplied the pretraining that made those detectors converge~\cite{deng2009imagenet}. COCO raised the instance count per image and became the standard mask benchmark~\cite{lin2014coco}. A warehouse of identical crates is a different distribution from COCO. The features transfer; the decision threshold and the color head do not, which is why a prototype head trained on crate colors is part of the update rather than a reuse of the COCO classifier.

\subsection{The one-stage mask and the YOLO line}
\label{sec:yolo}

Two-stage mask heads are accurate and, for a forklift that may clear the camera in under a second, often too sequential. YOLACT removed the region-wise mask head. It predicts a dictionary of prototype masks on the whole image and a vector of coefficients on each detection, then forms the instance mask as a linear combination of the prototypes, cropped to the box~\cite{bolya2019yolact}. YOLACT++ keeps that assembly and improves the prototypes and the backbone~\cite{bolya2022yolactpp}. The assembly is
\begin{equation}
M = \sigma\!\left( P\, c \right),
\label{eq:yolact}
\end{equation}
where $P$ is the stack of prototype masks, $c$ is the coefficient vector of one detection, and $\sigma$ is a sigmoid. The prototypes are shared. Individuality lives in $c$. FastSAM later showed that a YOLOv8 segmenter, which uses this same assembly, can be trained as a fast proposal generator for a segment-anything style task~\cite{zhao2023fastsam}.

Other one-stage designs are worth keeping in view because they expose different failure modes. SOLO predicts a mask on a spatial grid, so that location itself is the identity of the instance~\cite{wang2020solo}. SOLOv2 makes the mask weights dynamic and adds a matrix non-maximum suppression that is cheaper than a pixel-wise suppression~\cite{wang2020solov2}. CondInst generates the weights of a tiny mask head from the instance features, conditioned on the instance~\cite{tian2020condinst}. BlendMask combines a top-down attention on a bottom-up position-sensitive map~\cite{chen2020blendmask}. PolarMask represents the contour as rays from the center, which is attractive for compact blobs and brittle for a crate whose visible region is a rectangle with a bite taken out by occlusion~\cite{xie2020polarmask}. SparseInst activates a sparse set of instance queries and avoids non-maximum suppression~\cite{cheng2022sparseinst}. BoxInst shows that a projection loss and a pairwise color affinity can train a CondInst-style mask from boxes alone~\cite{tian2021boxinst}. The pairwise term is relevant here: plastic crates are nearly constant in color inside the face, so a color-affinity prior agrees with the object, and it also agrees with the failure mode in which two adjacent crates of the same color are glued into one mask.

The YOLO family itself is the practical backbone of the baseline used in this thesis. After the original grid regression~\cite{redmon2016yolo}, YOLO9000 introduced joint training and anchor boxes~\cite{redmon2017yolo9000}, and YOLOv3 moved to a feature pyramid and a stronger backbone~\cite{redmon2018yolov3}. YOLOv4 organized a large set of training and architecture choices that had been scattered across papers~\cite{bochkovskiy2020yolov4}. YOLOv5 made that line reproducible as a maintained implementation~\cite{jocher2020yolov5}. YOLOv6 reworked the neck and the loss for industrial deployment~\cite{li2022yolov6}. YOLOX replaced the coupled head and the anchors with a decoupled head and an anchor-free assignment~\cite{ge2021yolox}. YOLOv7 studied which of the so-called freebies survive re-parameterization and set a high bar for real-time detection~\cite{wang2023yolov7}. YOLOv8, released as software rather than as a single paper, is the version whose segmentation head follows~\eqref{eq:yolact} and whose detection head is decoupled~\cite{jocher2023yolov8}. YOLOv9 introduced programmable gradient information so that the loss sees a more complete view of the input, and showed a gain on the mask task as well as on boxes~\cite{wang2024yolov9}. YOLOv10 removed the need for non-maximum suppression by a consistent dual assignment~\cite{wang2024yolov10}. YOLO11 continues the same implementation line~\cite{jocher2024yolo11}. Real-time transformer detectors, from DETR~\cite{carion2020detr} through deformable attention~\cite{zhu2021deformable} and denoising queries~\cite{zhang2023dino} to RT-DETR~\cite{zhao2024rtdetr}, have closed the speed gap. They are not the baseline here, because the crate system needs a mask, a stable box center, and a small extra head, and the YOLACT-style segmenter already provides the first two.

The feature pyramid that these detectors share is the structure of Lin and colleagues~\cite{lin2017fpn}: a top-down path that carries semantics down to high-resolution maps, and lateral connections that restore detail. EfficientDet replaced the plain sum of those laterals with a weighted bidirectional pyramid~\cite{tan2020efficientdet}. CSPNet reduced duplicate gradient flow by splitting the feature map, and that split is the ancestor of the cross-stage modules in the later YOLO backbones~\cite{wang2020cspnet}. Channel attention of the squeeze-and-excitation kind~\cite{hu2018senet}, the convolutional block attention module~\cite{woo2018cbam}, and coordinate attention~\cite{hou2021coordatt} are the usual tools when a designer wants the network to prefer one spatial axis. The vertical-feature neck of YOLO-vf-seg, defined properly in Section~\ref{sec:problem-gap}, is a deliberate use of that last idea: on a side view the informative edges are horizontal, so the neck should be allowed to treat the vertical coordinate differently from the horizontal one.

Backbones have their own literature, and the warehouse papers cited in the next section borrow from it. Residual networks made the depth of the backbone a free parameter rather than an optimization accident~\cite{he2016resnet}. MobileNet and MobileNetV2 showed that depthwise separable convolution and inverted residuals keep accuracy when the computer is small~\cite{howard2017mobilenet,sandler2018mobilenetv2}. EfficientNet scaled depth, width, and resolution together~\cite{tan2019efficientnet}. ConvNeXt modernized a plain convolutional network with the training recipe of transformers~\cite{liu2022convnext}. The vision transformer and the Swin transformer replaced the sliding window with attention over patches~\cite{dosovitskiy2021vit,liu2021swin}, using the attention mechanism introduced for sequence transduction~\cite{vaswani2017attention}. RepVGG trains a multi-branch block and folds it into a single convolution at inference~\cite{ding2021repvgg}. That fold is why several recent warehouse models can quote a high frame rate after having trained a heavier block. The present design does not replace the backbone. It adds a color head whose parameter count is a fraction of one percent of a medium segmenter, and a counting stage that is not a network at all.

Box regression losses affect whether two stacked crates receive two centers or one merged box. Generalized IoU adds the empty area of the smallest enclosing box to the ordinary overlap penalty~\cite{rezatofighi2019giou}. Distance-IoU and the complete-IoU variant add the distance between centers and the aspect ratio, so that a non-overlapping prediction still receives a gradient~\cite{zheng2020diou}. Wise-IoU modulates that geometry by a focusing coefficient that depends on how the rest of the batch is behaving~\cite{tong2023wiou}. IoU-Net argued that the classification score is a poor proxy for localization quality and predicted the overlap directly~\cite{jiang2018iounet}. Generalized focal loss and VarifocalNet fold a localization quality into the classification target so that non-maximum suppression ranks boxes by a score that means what the suppression thinks it means~\cite{li2020gfl,zhang2021vfnet}. For crates, a high classification score on a box that straddles two stacked faces is exactly the error that produces a missed count. The lattice stage can repair some of those misses. It cannot repair them if the center it receives lies between two true cells.

Mask losses have a similar tension between region overlap and boundary fidelity. The Dice coefficient~\cite{dice1945} and its soft version in V-Net~\cite{milletari2016vnet} emphasize overlap and are stable when the crate occupies a modest part of the window. A generalized Dice reweights classes by the inverse of their volume~\cite{sudre2017gdice}. Overlap losses under-penalize a boundary that is off by a few pixels, and a few pixels on a crate lip are the difference between one instance and two. PointRend treats the uncertain boundary pixels as a rendering problem and classifies them at higher resolution~\cite{kirillov2020pointrend}. Boundary patch refinement crops small windows along the predicted contour and re-segments those windows~\cite{tang2021bpr}. Both are compatible with YOLO-vf-seg as a later refinement. They are not the updates proposed here, because the dominant counting errors in a regular stack are missed cells and wrong colors, not a contour that is cosmetically rough.

\subsection{Counting what is visible and what is hidden}
\label{sec:counting-principles}

The oldest learned counter regresses a density map whose integral is the count. Lempitsky and Zisserman introduced that idea so that a model could be trained from dotted annotations~\cite{lempitsky2010count}. A multi-column network estimates the density at several receptive fields when the objects vary in scale~\cite{zhang2016mcnn}. CSRNet keeps a single deeper column and uses dilated convolution to enlarge the field without losing resolution~\cite{li2018csrnet}. Later losses match the predicted and target distributions rather than only their pixel values, which reduces the tendency to place a blurry mass where a sharp peak should be~\cite{wang2020dmcount}. A composition loss combines counting, density, and localization~\cite{idrees2018composition}. Counting in the wild showed that the same idea extends past crowds when the annotations are still points~\cite{arteta2016counting}.

A density map is the wrong intermediate for typed crates. The integral of a density does not remember which color produced which mass of the integral, and a side view of a stack violates the assumption that every object contributes a visible dot. The bottom row is partly behind the forks. Crates in the second and third rank of the pallet are not in the image at all. Integrating a density over the visible face estimates the visible count, and only under the further assumption that occlusion eats a similar fraction of every crate. That assumption is false at the bottom and at the top of the frame.

Detection counting is the other extreme. Each surviving mask contributes one to a tally. It is simple, it preserves class, and it is biased low by every miss. Area-ratio counting, discussed with the warehouse systems in the next section, tries to recover partial crates by dividing a visible area by a nominal area. It recovers a fraction of a crate and then has to round. Geometry-aware counting uses the fact that the objects are the same size. That is the opening this thesis uses. The front face of a beer-crate stack is a lattice with known pitch. A missed detection is often a hole in an otherwise regular grid, and a hole with supported neighbors is a different event from an empty region of the aisle.

Hidden objects can also be counted from volume. Dumery and colleagues decompose the count of identical stacked objects into a reconstructed volume times an occupancy ratio estimated from a depth map, and they show that people are poor at the same task~\cite{dumery2024stacked}. Their method assumes a container seen from several views, objects that are approximately identical, and a depth estimator. A brewery forklift presents a different bargain. The objects are identical only within a color, the colors matter, the useful view is the side rather than the top, and the hidden dimension is a small integer: how many crates deep the pallet is. Multiplying an occupied front lattice by that integer is a coarser, and for this scene a more direct, use of the same idea that hidden extent should be measured rather than hallucinated pixel by pixel.

Associative embedding groups detections that belong to the same structured object by a learned tag~\cite{newell2017ae}. A Hungarian assignment matches two sets at minimum cost~\cite{kuhn1955hungarian}. PLC does not need either construction for the front face, because the lattice supplies the grouping. Both remain relevant if a later version of the system has to associate the same physical crate across a long passage with a moving camera. The estimate in this thesis uses a simpler rule, and Chapter~\ref{ch:system} reports that the simpler temporal rule over-counts.

\subsection{Color constancy, and why a uniform face breaks it}
\label{sec:color-principles}

Color constancy estimates $\mathbf{L}$ in~\eqref{eq:diagonal} so that a corrected image $\mathbf{I}(x)\oslash\hat{\mathbf{L}}$ depends on reflectance. The gray-world hypothesis says that the average reflectance in the scene is gray, so the average of $\mathbf{I}$ is a sufficient estimate of the illuminant~\cite{buchsbaum1980grayworld}. The Minkowski family generalizes that average. With $p\ge 1$,
\begin{equation}
\hat{\ell}_c
= \left(
\frac{1}{\lvert\Omega\rvert}
\sum_{x\in\Omega} \lvert I_c(x)\rvert^{p}
\right)^{1/p}.
\label{eq:minkowski}
\end{equation}
The case $p=1$ is gray-world. The case $p\to\infty$ is max-RGB, the channel-wise maximum. Shades-of-gray argues that a moderate $p$, commonly six, is a better default across photographs than either extreme~\cite{finlayson2004shades}. Edge-based variants apply the same norms to image derivatives so that large flat surfaces, which dominate a gray-world average and carry little information about the illuminant, do not drown the estimate~\cite{vandeweijer2007edge}. Retinex theory is the perceptual ancestor of these algorithms: lightness is inferred from spatial comparisons rather than from the absolute intensity at a point~\cite{land1971retinex}.

Learned estimators replace the norm. Oh and Kim cast illuminant estimation as classification over a discrete set of lights~\cite{oh2017colorclass}. FC4 predicts a local illuminant and a confidence at every patch and pools them with the confidence as a weight, precisely because many patches are consistent with almost any lamp~\cite{hu2017fc4}. Fast Fourier color constancy reduces illuminant estimation to localization on a torus and runs fast enough for a camera pipeline~\cite{barron2017ffcc}. The survey of Gijsenij, Gevers, and van de Weijer remains the right map of the assumptions behind the classical methods~\cite{gijsenij2011survey}. Afifi and Brown showed, from the other direction, that a network trained for recognition loses accuracy when the input white balance is wrong, so color constancy is not a cosmetic preprocessing step~\cite{afifi2019colorbias}.

The crate face inverts the assumption that makes~\eqref{eq:minkowski} valid. Gray-world and shades-of-gray need a mixture of reflectances. A beer crate is one reflectance, plus a logo that covers little of the face. On a masked face the sum in~\eqref{eq:minkowski} estimates $\mathbf{L}\odot\mathbf{R}$, not $\mathbf{L}$. Dividing $\mathbf{I}$ by that estimate pushes the corrected color toward gray and erases the type. This is not a failure of implementation. It is the hypothesis being false on the domain. A global estimate on the whole frame is not a repair: the frame contains a green crate, an amber crate, a gray floor, and a yellow safety line, and the average of those surfaces is not the illuminant on any one crate. Chapter~\ref{ch:methods} therefore keeps the Minkowski norm only as a feature channel, computed inside the mask, and trains a classifier that is allowed to use the residual chromaticity rather than being forced to trust the corrected color as if it were reflectance.

\subsection{Occlusion boundaries}
\label{sec:occlusion-principles}

When one crate overlaps another, the contour in the image is a mixture of the true object boundary and the occlusion boundary. Standard mask heads do not distinguish the two, so the occluded crate is trained to look like its visible fragment. BCNet models the local region as two overlapping layers and regresses the occluder and the occludee with separate graph-convolution branches, which decouples the two boundaries~\cite{ke2021bcnet}. MaskFormer and Mask2Former recast segmentation as mask classification with a transformer decoder, and the masked attention of Mask2Former restricts each query to the region it already believes is its own object~\cite{cheng2021maskformer,cheng2022mask2former}. Mask DINO carries the same decoder into a detector that shares queries between boxes and masks~\cite{zhang2023maskdino}. The segment-anything model and its video successor show that a large promptable segmenter can produce masks for objects it was not trained to name~\cite{kirillov2023sam,ravi2024sam2}. Naming remains our problem. A foundation mask is a plausible front end. It does not know that amber and yellow are different stock-keeping units, and it does not know that a hole in a regular stack is probably a missed crate rather than a hole in the load.

The side-view stack adds a regularity that generic occlusion models do not use. Crates of one product are the same size, they are placed on a pallet in a grid, and the hidden ones continue that grid in depth. A bilayer head can paint the visible fragment of an occluded crate more honestly. It cannot see the crate behind the front face. The count of those crates is a geometric statement, not a segmentation statement. That split, between a mask model for what is visible and a lattice model for the regular structure the mask only samples, is the design choice defended in Section~\ref{sec:problem-gap}.

\begin{figure}[t]
\centering
\begin{tikzpicture}[
  font=\small,
  box/.style={draw, rounded corners=2pt, align=center, inner sep=5pt, minimum height=8mm},
  arr/.style={-{Stealth}, thick}
]
  \node[box, fill=blue!8] (img) {Side-view image};
  \node[box, fill=blue!8, right=8mm of img] (seg) {Instance masks\\and centers};
  \node[box, fill=orange!15, right=8mm of seg] (col) {Color type\\per mask};
  \node[box, fill=green!12, below=8mm of seg] (lat) {Front-face lattice};
  \node[box, fill=green!12, right=8mm of lat] (dep) {Depth reading};
  \node[box, fill=black!6, below=8mm of col] (out) {Count per color\\and total count};
  \draw[arr] (img) -- (seg);
  \draw[arr] (seg) -- (col);
  \draw[arr] (seg) -- (lat);
  \draw[arr] (lat) -- (dep);
  \draw[arr] (col) -- (out);
  \draw[arr] (dep) -- (out);
\end{tikzpicture}
\caption{The computation this thesis commits to. Segmentation proposes visible crates. Color type is a second decision on each mask. The count is not the number of masks: it is the number of occupied cells on a rectified lattice, times the depth of the pallet. The two shaded branches are the updates developed in Chapter~\ref{ch:methods}.}
\label{fig:pipeline-overview}
\end{figure}

\section{Previous systems for cartons, warehouses, and stacked goods}
\label{sec:prior-systems}

\subsection{A dataset of stacked cartons}
\label{sec:scd}

Public detection sets barely contain a loaded pallet. Yang and colleagues built the Stacked Carton Dataset to close that gap: 16{,}136 images and about 250{,}000 instance masks, split between warehouse photographs and images gathered online, with labels that mark a carton as inner or outer and as occluded or fully visible~\cite{yang2022scd}. The accompanying detector starts from RetinaNet and adds two modules. Offset prediction between classification and localization corrects the score used by non-maximum suppression so that a well-localized box is not discarded in favor of a worse box with a higher class score. Boundary-guided supervision uses the mask contour as an extra training signal so that repeated printed textures do not glue neighboring cartons together. Both ideas transfer in spirit to beer crates. The class score of a YOLO head is not a count, and the seam between two amber crates is exactly the repeated-texture problem, only worse, because the texture is a flat color rather than a printed logo.

The benchmark on that dataset also records how standard detectors behave on stacks. Two-stage and one-stage detectors, DETR, YOLOX, and CenterNet are all evaluated, and so are several mask models. The spread of those results is a warning against quoting a COCO mask average precision as if it were a warehouse count accuracy. A model can gain several points of box average precision and still merge the seam that the count depends on. The carton dataset does not, however, label body color as the quantity to be inventoried, and it does not place the camera on the path of a forklift. It is the right evidence that stacked boxes are a segmentation problem. It is not yet the problem in Section~\ref{sec:scene}.

\subsection{Instance segmentation of a warehouse scene}
\label{sec:warehouse-yolo}

Zhang and Tan modify YOLOv8n for instance segmentation of an indoor warehouse~\cite{zhang2025warehouse}. They replace the feature block with a RepVGG-style block, insert a multi-scale convolutional attention between the backbone and the decoder, and train with Wise-IoU. On their CargoV10 set the mask mean average precision at an overlap of one half is reported as 70.5 percent, 5.1 points above the unmodified baseline, at 84.8 frames per second. The classes include goods, labels, forklifts, pallets, and people. That class list is the neighboring problem, not the same problem. A forklift in their images is an object to be segmented. In this thesis the forklift is the carrier of the object, and the camera is aimed at the load. Their result still matters for two engineering reasons. First, a small YOLO segmenter sustains a warehouse frame rate after the backbone has been re-parameterized, so a color head of a few thousand parameters will not be what breaks the timing budget. Second, Wise-IoU and an attention block aimed at adjacent similar objects are evidence that the community already treats ``neighboring instances that look alike'' as the warehouse failure mode. Beer crates of one color are that failure mode in its pure form.

CheapConv pushes the same industrial constraint from the other side: parcel segmentation that has to run in real time on a logistics line, with an architecture pared down for that budget~\cite{zou2025cheapconv}. The design lesson is identical. Accuracy that arrives only in a large transformer decoder does not automatically survive contact with a line that cannot wait. The updates in this thesis are shaped by that constraint. ICPH is a linear map on a pooled feature. PLC is arithmetic on a few dozen centers.

\subsection{Area ratios and three-dimensional occupancy}
\label{sec:area-and-3d}

Hu, Wang, and Yang count stacked goods by an area ratio~\cite{hu2025arbca}. An instance segmenter, in their experiments Mask2Former, produces a mask for each visible piece of the stack. The area of that mask is divided by the area the object would have presented had it been fully visible, and the ratios are accumulated into an inventory. On a tobacco-cargo stacking set they report an inventory accuracy of 93.90 percent. The method is aimed at occlusion and at dark warehouses, which are both present in a brewery aisle.

The area ratio is the natural baseline for PLC, and Chapter~\ref{ch:system} estimates it under the same observation model. Its weakness on a regular crate stack is structural. A mask that covers 0.86 of a face contributes 0.86 to the count only if the missing 0.14 is occlusion of a present crate. The same 0.14 is missing when the segmenter simply stopped at a shadow. The ratio then needs a rounding rule, and rounding errors add across a stack of twenty crates. A lattice does not round a fraction of a cell. It decides that a cell is occupied or not, using the pitch that the area ratio never consults. The area ratio remains the right tool when the goods are irregular and no pitch exists. Tobacco bales and mixed cartons can look like that. A single family of beer crates does not.

The three-dimensional counter of Dumery and colleagues is the other published answer to hidden objects~\cite{dumery2024stacked}. Volume times occupancy is powerful when the stack is a jumble inside a bin and several camera views can be taken. It is a poor fit when the hidden axis is an integer depth of identical crates, the camera sees one side during a passage of a second or two, and the operator already knows, or can read from the pallet pattern, whether the load is one, two, or three crates deep. PLC takes the integer and refuses the volumetric reconstruction. The comparison is not that one paper is wrong. The comparison is that the regularity of the brewery load is information, and a method that ignores it has to relearn the pitch from pixels on every frame.

\subsection{Rotated, lightweight detection of logistics boxes}
\label{sec:rew}

REW-YOLO detects logistics boxes with a lightweight rotated detector~\cite{wang2025rewyolo}. A RepViT-style block replaces part of the YOLOv8 convolution, an efficient local attention is built into the feature pyramid, a rotation angle is regressed, and Wise-IoU weights the box loss. On their BOX-data the model reaches 90.2 percent mean average precision at one half overlap, with 2.18 million parameters and about 131 frames per second. On the occluded and rotated subset they report 91.3 percent precision and 86.6 percent recall. The paper is explicit that YOLOv8 misses boxes whose occlusion exceeds roughly two thirds, because the edge that the detector needed has disappeared.

Rotation matters less for the front face of a pallet that was stacked square to the forks, and it matters immediately if the forklift yaws while it passes the camera. A homography estimated once per passage absorbs a modest yaw if the face stays planar. It does not absorb a crate that was loaded askew inside an otherwise square stack. The present design treats that as a residual error of the lattice snap, visible in Chapter~\ref{ch:system} as a center noise of a few centimetres, rather than as an angle to be regressed on every box. If a deployment sees large in-plane rotation, the REW-YOLO head is a compatible replacement for the box branch. It still would not name the color under a changing lamp, and it still would not multiply by depth. Those are the two holes the present updates are cut to fill.

REW-YOLO also documents the data augmentation that logistics papers now treat as mandatory: random rectangular occluders covering 5 to 25 percent of a box, and perturbations of brightness, contrast, hue, and saturation. The color estimate in this thesis is not trained on those augmentations of real photographs. It is trained inside a generative model that applies a random illuminant gain, a shadow, and a strip of a neighboring color. The augmentation list is evidence that illumination and occlusion are already recognized as the nuisance variables. The estimate asks how large the illuminant nuisance can grow before a chromatic prototype loses the type.

\subsection{What crowd counting and global white balance do not give the aisle}
\label{sec:negative}

It is worth stating the negative conclusions in one place, because they determine the architecture.

A crowd counter trained on point annotations will, at best, estimate how many crate-like blobs are visible. It will not emit a count per color. It will not add the rank of crates behind the front face. Dilated backbones and distribution-matching losses improve the visible tally on congested scenes~\cite{li2018csrnet,wang2020dmcount}. They do not add the missing state. Using one because it is the strongest public counter would be a category error.

A global white balance, however accurate on the benchmarks used by FC4 or fast Fourier color constancy~\cite{hu2017fc4,barron2017ffcc}, answers a different question: which single gain makes the whole photograph look neutral. The aisle does not need a neutral photograph. It needs the reflectance class of each crate. Applying a single gain can even hurt, if the gain is set by the floor and the safety paint rather than by the light on the stack. Per-mask features are the smallest domain on which the question is well posed.

A promptable segmenter can be pointed at the stack and will often return a plausible set of visible regions~\cite{kirillov2023sam,ravi2024sam2}. Without a color prototype and without a lattice, the regions are an unlabeled partition. The rest of the thesis is that partition's missing second and third stages.

\subsection{Where the published systems stop}
\label{sec:gap-table}

\begin{table}[t]
\centering
\caption{What published systems decide, and what the brewery passage still needs. A dash means the quantity is outside the stated task.}
\label{tab:gap}
\footnotesize
\setlength{\tabcolsep}{3.5pt}
\begin{tabular}{@{}lcccc@{}}
\toprule
System & Visible instances & Color type & Hidden depth & Side-view lattice \\
\midrule
Mask R-CNN~\cite{he2017maskrcnn} & yes & only if trained & -- & -- \\
YOLACT / YOLO-seg~\cite{bolya2019yolact,jocher2023yolov8} & yes & only if trained & -- & -- \\
Mask2Former~\cite{cheng2022mask2former} & yes & only if trained & -- & -- \\
SCD carton detector~\cite{yang2022scd} & yes & -- & -- & boundary cue \\
Warehouse YOLOv8-seg~\cite{zhang2025warehouse} & yes & -- & -- & -- \\
Area-ratio inventory~\cite{hu2025arbca} & yes & -- & via area & -- \\
3D occupancy count~\cite{dumery2024stacked} & yes & -- & via volume & multi-view \\
REW-YOLO~\cite{wang2025rewyolo} & boxes & -- & -- & rotated box \\
Crowd density~\cite{li2018csrnet} & a scalar & -- & -- & -- \\
Global color constancy~\cite{barron2017ffcc} & -- & a global gain & -- & -- \\
\midrule
This thesis & yes & per mask & integer depth & front face \\
\bottomrule
\end{tabular}
\end{table}

\Cref{tab:gap} is the map of the related work used to choose the two updates. No listed system both names the crate by body color under a changing illuminant and turns a side view of a regular stack into a count that includes a depth integer. Several systems solve a neighboring piece extremely well. The design in Chapter~\ref{ch:methods} is intentionally a pair of small additions to a segmenter that already solves the visible-instance piece, rather than a new detector.

\subsection{Measures this thesis inherits, and measures it refuses}
\label{sec:metrics}

A detector paper reports a mean average precision. A counter reports an absolute error. The two numbers answer different questions, and the warehouse literature cited above is easy to misread when they are placed side by side. This subsection fixes the definitions that later chapters either use or explicitly decline.

Intersection over union of a predicted box $B$ and a ground-truth box $G$ is
\begin{equation}
\operatorname{IoU}(B,G)
= \frac{\lvert B\cap G\rvert}{\lvert B\cup G\rvert}.
\label{eq:iou}
\end{equation}
A detection is a true positive at threshold $\tau$ when its class is correct and $\operatorname{IoU}\ge\tau$. Precision is the fraction of predictions that are true positives. Recall is the fraction of ground-truth objects that are matched. Average precision integrates precision over recall for one class, and mean average precision averages that integral over classes~\cite{lin2014coco,everingham2010pascal}. Mask average precision replaces the box overlap by the mask overlap and is the number quoted for YOLACT and for the warehouse segmenter of Zhang and Tan~\cite{bolya2019yolact,zhang2025warehouse}. At $\tau=0.5$ the metric forgives a boundary that is off by a wide margin. At $\tau=0.75$ it begins to punish the seam error that, for two stacked crates, is the difference between a count of one and a count of two. A paper that reports only the looser threshold can improve while the tally stays wrong.

None of those averages is the objective in \Cref{def:task}. The tally is an integer per passage. The errors used in Chapter~\ref{ch:system} are therefore the errors of an integer. With $T$ trials, prediction $\hat{N}^{(t)}$ and truth $N^{(t)}$,
\begin{equation}
\begin{aligned}
\mathrm{MAE}
&= \frac{1}{T}\sum_{t=1}^{T}\lvert\hat{N}^{(t)}-N^{(t)}\rvert,\\
\mathrm{RMSE}
&= \left(\frac{1}{T}\sum_{t=1}^{T}\big(\hat{N}^{(t)}-N^{(t)}\big)^{2}\right)^{1/2},\\
\mathrm{Bias}
&= \frac{1}{T}\sum_{t=1}^{T}\big(\hat{N}^{(t)}-N^{(t)}\big),\\
\mathrm{MAPE}
&= \frac{1}{T}\sum_{t=1}^{T}\frac{\lvert\hat{N}^{(t)}-N^{(t)}\rvert}{N^{(t)}}.
\end{aligned}
\label{eq:count-metrics}
\end{equation}
The exact-match rate is the fraction of trials with $\hat{N}^{(t)}=N^{(t)}$. The ninetieth percentile of $\lvert\hat{N}-N\rvert$ describes the bad passage rather than the average passage. Bias is kept, despite being a smaller number than the mean absolute error whenever errors cancel, because a brewery that is short three crates and a brewery that is long three crates do not have the same problem. Density-map papers often report only a mean absolute error and a root mean square error on the crowd count~\cite{li2018csrnet,wang2020dmcount}. They can omit bias when the count is a single class of people. They should not omit it when the sign distinguishes hidden depth from invented cells.

Color uses top-1 accuracy, the fraction of masks whose predicted prototype equals the reflectance class, and a confusion matrix. Accuracy without the matrix is how amber and yellow disappear into a mean. The matrix is reported at the harshest illuminant for that reason. A mean average precision on color is not computed, because the color head is not a detector. It is a classifier on a mask the detector has already proposed. Mixing the two into one average would let a gain in box overlap hide a swap of stock-keeping units.

Two further metrics appear in the segmentation literature and are not the acceptance test here. The Dice coefficient between a mask and a ground-truth region~\cite{dice1945,milletari2016vnet} is a training loss and a boundary diagnostic. A high Dice on a merged pair of crates is a failure of counting that looks like a success of overlap, because the merged mask covers almost all of the plastic. Boundary F-scores of the kind used to argue for PointRend and for boundary patch refinement~\cite{kirillov2020pointrend,tang2021bpr} would be the right metric for a future refinement of the lip. They are not estimated, because this thesis does not claim a new boundary head.

The split is the one \Cref{tab:gap} already drew. Segmentation metrics judge YOLO-vf-seg as a proposer of regions. Counting metrics judge PLC as a proposer of integers. Classification metrics judge ICPH as a proposer of names. A single headline number would require a filmed passage in which all three are scored on the same crates. That passage is outside the estimate. The definitions above are how the estimate stays comparable to the prior work without pretending to have repeated it.

\begin{table}[t]
\centering
\caption{Which published number is allowed to stand in for which decision. The right-hand column is the quantity Chapter~\ref{ch:system} actually computes.}
\label{tab:metrics-map}
\small
\begin{tabular}{@{}p{0.34\textwidth}p{0.28\textwidth}p{0.28\textwidth}@{}}
\toprule
Published number & What it decides & What this thesis uses instead \\
\midrule
Box or mask mAP at $\tau=0.5$~\cite{lin2014coco} & Ranking of detectors on varied objects & Not computed. A seam error can hide inside $\tau=0.5$. \\
Inventory accuracy of an area ratio~\cite{hu2025arbca} & Fraction of stacks whose integer matches & Exact-match rate, plus MAE and bias \\
mAP of a rotated box detector~\cite{wang2025rewyolo} & Localization under yaw & Origin residual in centimetres \\
Angular illuminant error~\cite{barron2017ffcc,hu2017fc4} & Quality of a global gain & Top-1 type accuracy per mask \\
\bottomrule
\end{tabular}
\end{table}


\section{The problem this thesis addresses}
\label{sec:scene}

\subsection{The passage into the warehouse}
\label{sec:scene-geom}

At the end of the packaging line the product is no longer a bottle. It is a crate, and the crates are no longer on a conveyor. They are stacked on a pallet, the pallet is on the forks of a truck, and the truck is driving into the warehouse. That passage is the last moment at which the load is presented, face-on and in motion, to a camera that the building can own. A count taken later, from the aisle, sees the long side of a row of pallets and a great deal of occlusion between pallets. A count taken earlier, on the conveyor, sees single crates and does not see the stacking errors introduced when the pallet is built. The passage is the right place to close the tally.

The camera looks at the side of the load, approximately normal to the front face, as drawn on the title page and, with the quantities used later, in \Cref{fig:scene}. The front face is a grid. In the nominal geometry each cell is $0.40\,\mathrm{m}$ wide and $0.28\,\mathrm{m}$ tall, and each crate extends $0.30\,\mathrm{m}$ in the depth direction, away from the camera. A typical load is two to six crates tall, two to four crates wide on the face the camera sees, and one to three crates deep. The face is what the image contains. The depth is what the image does not contain. Color is a property of the plastic body: amber, green, red, blue, yellow, or white in the design studied here. Two colors on one pallet are normal, because a forklift may carry a mixed recovery of empty crates or a mixed dispatch. The system is required to count each color, not only the total.

\begin{figure}[t]
\centering
\begin{tikzpicture}[x=1cm,y=1cm,font=\small]
  \draw[thick] (0,0) -- (13.2,0) node[right] {\footnotesize floor};
  % forklift
  \draw[fill=black!10] (0.4,0.35) rectangle (3.1,1.7);
  \draw[fill=black!22] (0.9,1.7) rectangle (2.3,2.7);
  \draw[line width=1.1pt] (3.1,0.5) -- (4.7,0.5);
  \draw[line width=1.1pt] (3.1,0.75) -- (4.7,0.75);
  \node at (1.75,1.0) {\footnotesize truck};
  % pallet and stack
  \draw[fill=brown!50] (4.7,0.28) rectangle (9.4,0.85);
  \node at (7.05,0.56) {\footnotesize pallet};
  \foreach \r in {0,1,2,3} {
    \foreach \c in {0,1} {
      \pgfmathsetmacro{\yy}{0.95+\r*0.95}
      \pgfmathsetmacro{\xx}{4.95+\c*2.1}
      \draw[fill=orange!30] (\xx,\yy) rectangle ++(1.95,0.82);
    }
  }
  \draw[fill=green!28] (4.95,4.75) rectangle ++(1.95,0.82);
  \draw[fill=red!28] (7.05,4.75) rectangle ++(1.95,0.82);
  \node[align=center] at (7.05,5.85) {\footnotesize front face\\[-1pt]\footnotesize (visible)};
  % depth hint
  \draw[dashed] (9.4,0.85) -- (10.15,1.15) -- (10.15,5.55) -- (9.4,5.57);
  \draw[dashed] (9.4,0.85) -- (10.15,1.15);
  \node[align=left] at (11.35,3.2) {\footnotesize depth\\[-1pt]\footnotesize not in\\[-1pt]\footnotesize the image};
  % camera
  \draw[fill=black!80] (12.15,2.55) rectangle (13.05,3.05);
  \draw[-{Stealth}, thick] (12.15,2.8) -- (10.5,2.8);
  \node[align=center] at (12.6,3.45) {\footnotesize camera};
  \draw[{Stealth}-{Stealth}] (4.95,0.95) -- (4.95,5.57);
  \node[rotate=90, anchor=south] at (4.55,3.2) {\footnotesize height of stack};
\end{tikzpicture}
\caption{Side view of the passage. The camera sees the front face only. Crates behind that face, indicated by the dashed extension, change the count and do not change the image. The bottom of the face sits just above the forks and the pallet; the top of the face may leave the frame when the stack is tall.}
\label{fig:scene}
\end{figure}

The physical nuisances follow from that drawing. The forks and the pallet shadow the bottom row, so the bottom masks are shorter and darker than a textbook crate. The top row is the first to be clipped by the image border and the first to be grazed by motion blur if the shutter is slow relative to the truck. Crates in a column touch. Where two crates share a color, the only evidence of the seam is a thin lip, a change in shading, and the regular spacing of the lips above and below. Mixed colors make the seam easier and the color decision harder, because each mask is contaminated by a strip of the neighbor when the boundary is off by a few pixels. The lamp color is not under the designer's control. Dock doors contribute daylight. The interior contributes whatever lamp the warehouse installed. The truck may stop, or it may pass at a walking pace. The system has to produce a count for that one passage, not a count averaged over a day of well-lit stills.

\subsection{Formal task}
\label{sec:formal}

\begin{definition}[Passage count]
\label{def:task}
A passage supplies a short image sequence $I_{1},\ldots,I_{F}$, the intrinsics $K$, a homography $H$ from the image plane onto the front face of the stack, the nominal face size $(w,h)$, and a depth reading $\hat{d}\in\{1,2,3\}$. The load contains crates of classes $\mathcal{K}$. The required output is a count $\hat{N}_{k}$ for each $k\in\mathcal{K}$ and a total $\hat{N}=\sum_{k}\hat{N}_{k}$. A count is exact when $\hat{N}_{k}=N_{k}$ for every $k$. A total is exact when $\hat{N}=N$.
\end{definition}

The definition separates two accuracies on purpose. A swap of amber and yellow preserves the total and destroys the inventory. A missed cell destroys both. The estimate in Chapter~\ref{ch:system} measures color accuracy and total-count accuracy in separate trials, because the generative model draws type labels and geometric misses independently. A deployment would run them on the same masks. The definition is still the right target: both numbers are part of the product, and a system that reports only $\hat{N}$ has not finished the job.

The depth reading is part of the input, not part of the image. In the building it may come from the known pallet pattern, from a second, coarse sensor, or from the width of the load as seen by a top camera that this thesis does not design. The observation model used later is honest about that reading: it is correct with probability $0.90$ and otherwise off by one. The image system is not allowed to pretend that depth is free.

\begin{assumption}[Regular front face]
\label{ass:lattice}
Within one passage the front-face pitch is the nominal $(w,h)$, the face is close enough to a plane that one homography rectifies it, and a crate center that is detected falls nearer to its own cell than to a neighboring cell after rectification.
\end{assumption}

\begin{assumption}[Stable body color]
\label{ass:color}
Each class $k$ has a reflectance prototype. Illumination, shading, and a limited amount of boundary mixing change the observed color. They do not change the prototype that the class is supposed to match.
\end{assumption}

Both assumptions are design assumptions. They are true of a single family of returnable crates, stacked by a palletizer, photographed from a fixed gate. They are false of mixed cardboard of unknown size, and this thesis does not claim that case. \Cref{ass:lattice} is what allows PLC to snap centers to integers. \Cref{ass:color} is what allows ICPH to keep six prototypes rather than a continuous color regression.

\subsection{Requirements}
\label{sec:requirements}

The requirements below are the filter applied to every design choice in Chapter~\ref{ch:methods} and to the system sketch in Chapter~\ref{ch:system}.

\begin{enumerate}
\item \textbf{Type and total.} The output is a count per body color and the sum of those counts, as in \Cref{def:task}.
\item \textbf{One passage, one answer.} The truck does not return. The system may use a handful of frames while the face is in view. It may not average over the next truck.
\item \textbf{Side view only.} No requirement is stated for a top camera or a second side. Depth enters as a reading with a stated error rate.
\item \textbf{Illumination is not controlled.} A method that needs a gray card on every pallet, or a lamp of known spectrum, fails the requirement. A method that uses a calibration of the camera once per shift is acceptable.
\item \textbf{The seam is part of the task.} Two adjacent crates of the same color must be able to become two cells. A single mask over both is a counting error even if the pixels of plastic are correctly labeled as crate.
\item \textbf{The budget is a line, not a benchmark.} The segmenter must leave room for a color head and a counting stage inside the time the face is in view. A head that adds a few thousand parameters meets that budget. A second large network does not, unless a later measurement says there is spare time.
\item \textbf{Errors must have a sign that an operator can understand.} Under-counting hides product. Over-counting creates product. The estimator reports bias, not only an absolute error, so that a positive and a negative failure stay distinct.
\end{enumerate}

Requirement~7 is why Chapter~\ref{ch:system} reports signed bias next to every mean absolute error. A method that misses three crates and invents three others has zero bias and a disastrous inventory. Absolute error catches it. Bias catches a method that is systematically short, which is what raw mask counting does on a stack that is three crates deep.

\subsection{Where an unmodified segmenter fails}
\label{sec:fail}

An unmodified YOLO-style segmenter, trained to emit masks and a class, fails this passage in four repeatable ways.

The first failure is merged seams. The prototype assembly~\eqref{eq:yolact} is excellent at compact objects on COCO and weak at a long, low-contrast horizontal edge between two identical textures. If the coefficients of two stacked crates are similar and non-maximum suppression keeps one box, the tally is short by one for every such merge. Boundary refinement~\cite{kirillov2020pointrend,tang2021bpr} and occlusion-aware bilayer heads~\cite{ke2021bcnet} reduce this failure and do not use the fact that the lips are periodic. The vertical-feature neck defined below is the baseline's attempt to prefer those lips. It is not, by itself, a counter.

The second failure is color under the wrong lamp. A classifier trained on RGB crops of crates in one aisle learns $\mathbf{L}\odot\mathbf{R}$ for that aisle's $\mathbf{L}$. \Cref{eq:diagonal} says the decision boundary moves when the truck uses a different door. Global constancy does not put the boundary back, for the reason in Section~\ref{sec:color-principles}. The color error is not uniform across classes. Amber and yellow share a high red channel and a low blue channel. White sits near the neutral axis, so a small residual gain pushes it toward whichever chromatic class the noise favors. Any head that claims to solve type has to be estimated on those pairs, not only on a mean accuracy.

The third failure is the invisible rank. A perfect mask of the front face, multiplied by nothing, is the front count. If the pallet is three crates deep, that perfect mask is wrong by a factor of three. No improvement in mask average precision repairs a factor of three. This is the failure that makes a pure segmentation paper an incomplete answer to \Cref{def:task}, and it is the reason PLC multiplies by $\hat{d}$.

The fourth failure is the missed cell that is not a merged seam: the bottom crate in fork shadow, the top crate cut by the frame, the crate whose center was suppressed. These misses are correlated with position. A bottom row is harder than a middle row. A blurry frame is harder than a sharp one. A counter that treats every miss as an independent coin flip will be surprised by a bottom row that disappears together. The observation model used for the estimate makes that position dependence explicit, so the lattice is tested against the failure it is supposed to repair.

\begin{figure}[t]
\centering
\begin{tikzpicture}[font=\footnotesize, cell/.style={draw, minimum width=1.15cm, minimum height=0.78cm, inner sep=0pt}]
  % panel A merged
  \begin{scope}[shift={(0,0)}]
    \node at (1.7,3.35) {\small (a) merged seam};
    \node[cell, fill=orange!40, minimum height=1.7cm] at (1.15,1.7) {};
    \node[cell, fill=orange!25] at (2.4,2.15) {};
    \node[cell, fill=orange!25] at (2.4,1.25) {};
    \node[align=center] at (1.7,0.35) {one mask\\where two crates touch};
  \end{scope}
  % panel B color
  \begin{scope}[shift={(4.6,0)}]
    \node at (1.3,3.35) {\small (b) lamp shift};
    \node[cell, fill=yellow!55] at (0.7,2.15) {};
    \node[cell, fill=orange!60] at (1.95,2.15) {};
    \node[align=center] at (1.3,1.35) {amber and yellow\\move together};
    \node[align=center] at (1.3,0.45) {white drifts off\\the neutral axis};
  \end{scope}
  % panel C depth
  \begin{scope}[shift={(8.7,0)}]
    \node at (1.5,3.35) {\small (c) hidden depth};
    \foreach \i in {0,1,2} {
      \node[cell, fill=blue!20] at (0.7,0.7+\i*0.85) {};
    }
    \draw[dashed, fill=blue!8] (1.55,0.35) rectangle (2.7,3.15);
    \node[align=center] at (1.5,0.05) {front masks\\omit the rear ranks};
  \end{scope}
\end{tikzpicture}
\caption{Three of the four failures that define the problem. A merged seam shortens the tally by a whole crate. A lamp shift confuses the pairs whose reflectances are nearest. Hidden depth is invisible to any mask model that sees only the side face. The fourth failure, a shadowed or clipped cell, is the one the lattice gap rule is drawn to fill.}
\label{fig:failures}
\end{figure}

\subsection{YOLO-vf-seg, and the two updates}
\label{sec:problem-gap}

YOLO-vf-seg is the baseline this thesis updates. It is a one-stage instance segmenter of the YOLOv8-seg type: a cross-stage backbone, a feature pyramid, a decoupled detection head, and a prototype mask assembly of the form~\eqref{eq:yolact}. The letters ``vf'' name a vertical-feature gate in the neck. On a side view the signal that separates stacked crates is a set of nearly horizontal lips, periodic in the vertical image coordinate. The gate pools each feature map along the horizontal axis, passes the resulting column through a small convolution, and multiplies the map by a sigmoid of that column, expanded back across the width. In equations, for a feature map $U\in\mathbb{R}^{C\times H\times W}$,
\begin{equation}
a = \sigma\!\left( g( \operatorname{GAP}_{w}(U) ) \right) \in (0,1)^{C\times H},
\qquad
\tilde{U}_{c,i,j} = a_{c,i}\, U_{c,i,j},
\label{eq:vf}
\end{equation}
where $\operatorname{GAP}_{w}$ is the mean over the width and $g$ is a $1\times 3$ convolution along the vertical axis followed by a pointwise map back to $C$ channels. This is coordinate attention restricted to the axis that carries the stack~\cite{hou2021coordatt}, placed in the neck rather than used as a generic module. It does not classify color and it does not count. It makes the lips easier for the mask head to see. The rest of the thesis treats YOLO-vf-seg as given, and does not claim it as a new architecture.

The updated system adds two modules, drawn in \Cref{fig:pipeline-overview}.

The illumination-invariant chromatic prototype head, ICPH, reads the pixels inside each mask and the mask-pooled features of the neck. It builds the eight-dimensional chromatic feature of Chapter~\ref{ch:methods}, standardizes it, and scores six color prototypes with a linear classifier. The head exists because of the second failure mode above. It is small enough that requirement~6 is not in doubt: the parameter count estimated in Chapter~\ref{ch:system} is under ten thousand weights on top of a medium segmenter.

Perspective lattice completion, PLC, reads the mask centers after they have been mapped by $H$ into metres on the front face. It estimates the origin of the lattice, snaps each center to an integer cell, fills a hole that has at least two orthogonal neighbors inside the bounding rectangle of the detections, and multiplies the number of occupied cells by the depth reading. The module exists because of the third and fourth failure modes. It is not a network, and it has no training iteration. Its only geometric parameters are the nominal pitch $(w,h)$, which \Cref{ass:lattice} fixes.

A third idea was considered and set aside. Training the segmenter with an auxiliary density loss, so that the integral of a density agrees with the number of masks, repairs inconsistency between two visible-count estimators. It does not see behind the front face, and it does not use the pitch. A fourth idea, a second camera above the pallet, would see the depth directly and would violate requirement~3. Both ideas remain available if the assumptions fail in a particular building. They are not the contributions.

The contributions are deliberately unequal in kind. ICPH is a learned head, and the estimate of it is an estimate of a classifier under a color model. PLC is a geometric algorithm, and the estimate of it is an estimate of a counter under a detection model. Putting them in one network and quoting a single average precision would hide the separate ways they fail. Chapter~\ref{ch:methods} derives each one. Chapter~\ref{ch:system} builds the passage system around both, and then estimates them separately under models that are written down in full.

\subsection{What is in scope, and how the thesis is organized}
\label{sec:scope}

The thesis is a system design and an estimate of that design. It specifies the camera's relation to the load, the segmenter, the color head, the lattice, the depth reading, and the way a passage is reduced to one answer. It does not report a field trial. The numbers in Chapter~\ref{ch:system} are produced by the Monte Carlo instrument \texttt{simulate\_estimation.py} under the observation model written there, with the seed fixed. They are estimates of the two modules. They are not accuracies measured on a brewery floor, and the text does not call them that.

Several related problems are named so that they are not silently included. Bottle-level inspection inside the crate, reading of printed lot codes, detection of a broken crate, and identity of the forklift are out of scope. So is any claim about a cardboard carton of unknown size. The lattice pitch is an input. If a second family of crates with a different pitch shares the aisle, the system needs a pitch hypothesis per family; that extension is described and not estimated.

The organization follows the three-chapter limit. This chapter has set out the principles, the prior systems, and the problem. \Cref{ch:methods} derives ICPH and then PLC, with the losses, the algorithms, and the complexity. \Cref{ch:system} specifies the passage system that holds the two modules, and then reports the estimate: protocol, tables, the worked passage, the parameter cost, and the limits of what the estimate can support.

The standard that the later chapters are written to is the standard of this problem statement. A module earns its place if it attacks one of the four failures in \Cref{fig:failures} or the invisible rank behind them, if it respects the side-view and timing requirements, and if its errors can be shown with a sign. ICPH and PLC are the two modules that meet that standard. The next chapter writes them down.

\chapter{Two Updates of YOLO-vf-seg}
\label{ch:methods}

The baseline defined in Section~\ref{sec:problem-gap} already proposes a mask and a center for each visible crate. This chapter adds the two computations that turn those proposals into a typed inventory. Section~\ref{sec:icph} derives the illumination-invariant chromatic prototype head. Section~\ref{sec:plc} derives perspective lattice completion. Each section states the quantity being computed, the reason a simpler quantity is the wrong one, the algorithm, the cost, and the way the module fails when its assumption fails. The numerical estimate is withheld until Chapter~\ref{ch:system}, where the observation model can be written next to the tables. A formula that appears in both places has the same meaning in both places.

\section{Illumination-invariant chromatic prototype head}
\label{sec:icph}

\subsection{Where the two modules sit on the forward pass}

Before either update is derived, the forward pass they attach to should be fixed in dimensions, so that ``a small head'' and ``a center'' are quantities rather than slogans. The pass described here is the YOLOv8-seg assembly with the vertical-feature gate of~\eqref{eq:vf}. It is the baseline. Nothing in this subsection is claimed as a contribution. The contributions begin at Section~\ref{sec:icph}.

An input frame is resized to the network's square working resolution, denoted $S\times S$ with three channels. The value of $S$ is a deployment choice, commonly $640$. The footprint calculation in Section~\ref{sec:footprint} shows that the lattice is not starved of pixels at the gate, so the reason to keep $S$ large is the lip between two stacked crates, not the metre-scale pitch. The backbone, a cross-stage convolutional network of the kind introduced for efficient gradient flow~\cite{wang2020cspnet} and used by the maintained YOLO implementations~\cite{jocher2023yolov8}, emits three maps whose strides are $8$, $16$, and $32$ pixels. On a $640$ input those maps are $80\times 80$, $40\times 40$, and $20\times 20$. The neck builds the feature pyramid~\cite{lin2017fpn} and, on each of the three maps, applies the vertical gate
\begin{equation}
a = \sigma\!\left( g(\operatorname{GAP}_{w}(U)) \right),
\qquad
\tilde U_{c,i,j} = a_{c,i}\, U_{c,i,j}.
\label{eq:vf-again}
\end{equation}
The pooled column $\operatorname{GAP}_{w}(U)$ has shape $C\times H\times 1$. A convolution of kernel $1\times 3$ along $H$ mixes each row with the lip above it and the lip below it, which is the only place in the baseline where the stack's periodicity is written down as an architectural bias. The sigmoid keeps the gate in $(0,1)$ so that a row can be suppressed and cannot be amplified without bound. A row that contains a horizontal lip receives a larger $a$ if $g$ learns to respond to the vertical gradient of the pooled column. A row that contains only the flat plastic of a crate face receives a smaller $a$. The gate does not know the pitch in metres. It knows only that the informative variation, on this camera, runs down the image.

The detection head is decoupled. At each spatial site and at each stride it predicts an objectness, a class distribution over whatever classes the segmenter was trained to emit, and a box. For the brewery the segmenter's class may be the single class ``crate.'' Color is not this distribution. Color is ICPH, precisely so that a change of lamp does not have to be solved by the same softmax that decides whether a site is a crate. The box is parameterized as an offset from the site center plus a width and a height, the anchor-free form used since YOLOX~\cite{ge2021yolox}. Non-maximum suppression then keeps one box when two sites describe the same crate. The surviving box is the window inside which the mask is cropped, and its center is the fallback center if the mask is empty. The mask head, in parallel, predicts $K_{p}=32$ prototype masks at stride $4$, so on a $640$ input the prototypes are $160\times 160$, and one coefficient vector $c\in\mathbb{R}^{32}$ per detection. The mask is~\eqref{eq:yolact}, cropped to the box and thresholded at one half. The center passed to PLC is the first moment of the thresholded mask,
\begin{equation}
\bar u = \frac{\sum_{x\in M} u(x)}{\lvert M\rvert},
\qquad
\bar v = \frac{\sum_{x\in M} v(x)}{\lvert M\rvert},
\label{eq:moment}
\end{equation}
then mapped by $H$. Using the moment rather than the box center lets a crate that is half in shadow pull its reported center toward the visible plastic, which is usually toward the interior of the stack rather than toward the aisle. That pull is helpful when the shadow is on the outside of the face and harmful when the shadow is the seam, because the moment then slides toward the neighbor. The $3.5\,\mathrm{cm}$ noise in the counting model is a stand-in for both effects. It is not a measurement of~\eqref{eq:moment} on real masks.

ICPH reads $M$ and, in the full specification, the mask-pooled vector~\eqref{eq:pool} from the finest neck map. It does not read the prototype coefficients. Those coefficients are a spatial code, shared in meaning with YOLACT's prototypes~\cite{bolya2019yolact}, and they are not a color. Reusing $c$ as a color feature would entangle the seam with the lamp. The head's output is a class and a margin, attached to the same index as the center, so that cell $(r,c)$ and class $\hat k$ refer to one detection until the gap rule invents a cell that has no detection. Invented cells are the only ones that inherit a class, by the neighbor rule in~\eqref{eq:per-color}.

\begin{figure}[t]
\centering
\begin{tikzpicture}[
  font=\footnotesize,
  box/.style={draw, rounded corners=1.5pt, align=center, inner sep=3pt, minimum height=8mm, minimum width=18mm},
  arr/.style={-{Stealth}, semithick}
]
  \node[box, fill=black!6] (in) {Image\\$S\times S\times 3$};
  \node[box, fill=black!6, right=4mm of in] (bb) {Backbone\\strides 8, 16, 32};
  \node[box, fill=blue!10, right=4mm of bb] (vf) {Vertical gate\\on each level};
  \node[box, fill=black!6, right=4mm of vf] (det) {Decoupled\\boxes};
  \node[box, fill=black!6, below=7mm of det] (msk) {32 prototypes\\and coefficients};
  \node[box, fill=orange!15, below=7mm of vf] (icph) {ICPH\\class, margin};
  \node[box, fill=green!12, below=7mm of bb] (plc) {PLC\\cells $\times$ depth};
  \draw[arr] (in) -- (bb);
  \draw[arr] (bb) -- (vf);
  \draw[arr] (vf) -- (det);
  \draw[arr] (det) -- (msk);
  \draw[arr] (msk) -- (icph);
  \draw[arr] (msk.west) -- ++(-2mm,0) |- (plc.east);
  \node[align=center, font=\scriptsize] at ($(msk.west)+(-1.15,-0.55)$) {centers};
  \node[align=center, font=\scriptsize] at ($(msk.south)+(0,-0.45)$) {pixels in $M$};
\end{tikzpicture}
\caption{Forward pass of the updated model. The backbone and the vertical gate are the baseline YOLO-vf-seg. The prototype assembly produces the mask. ICPH and PLC both read that mask and do not read each other. Depth enters PLC from outside the figure.}
\label{fig:forward}
\end{figure}

The cost of being this explicit is that several widely discussed alternatives are now visibly in the wrong place. A second class head on the detection branch would name the color from the box features, including background inside the box, and would be trained by the same assignment that trains objectness. ICPH is downstream of the mask so that the pixels are the crate. A density head on the pyramid would estimate a visible count and would never see $\hat d$. PLC is downstream of the centers so that the pitch can be applied in metres. A global constancy preprocessing in front of the backbone would rescale the whole frame before the gate saw it, and \Cref{prop:uniform} says a per-face rescaling is already destructive if it is used as the classifier. No constancy step sits in front of the backbone. The camera gain is locked, and the only place an illuminant is estimated is inside $M$, as one block of $\mathbf{f}$.

The numbers that a deployer tunes, and that this thesis does not claim to have tuned on images, are $S$, the non-maximum-suppression overlap, the mask threshold of one half, and $K_{p}$. Moving $K_{p}$ from $32$ to $16$ halves the prototype tensor and, on COCO-like masks, usually costs a small amount of boundary detail~\cite{bolya2019yolact,zhao2023fastsam}. On a crate face the boundary that matters is one horizontal lip. Whether $16$ prototypes can express that lip is a question for a labeled gate, not for the Monte Carlo model, which never renders a lip. The model begins at the center. The architecture above is what has to produce a center that \Cref{prop:snap} can trust.


\subsection{The decision the head is asked to make}
\label{sec:icph-decision}

Let $M\subset\Omega$ be the set of pixels that the segmenter has marked as one crate, and let $\mathbf{I}(x)$ be the observed color at those pixels. ICPH returns a class $\hat{k}\in\mathcal{K}$ with $\lvert\mathcal{K}\rvert=6$, the six body colors of \Cref{ass:color}. It does not return a corrected image, and it does not return an illuminant for the whole frame. The inventory needs a name per mask. Anything else the head computes is an intermediate that exists to make that name stable.

The head is attached where the mask already exists, for a reason that is easy to understate. A color decision taken on a box, rather than on a mask, mixes the crate with the background that the box contains: the gap between slats, a piece of the crate behind, the forks at the bottom row. A color decision taken on the whole frame mixes every crate on the pallet. The mask is the smallest region for which \Cref{ass:color} is even approximately true. ICPH therefore refuses to run on a detection that has no mask. If the segmenter emits only a box, the system falls back to the raw prototype match described below and marks the type as low confidence. That fallback is a degradation, not the design.

\begin{figure}[t]
\centering
\begin{tikzpicture}[
  font=\small,
  box/.style={draw, rounded corners=2pt, align=center, inner sep=4pt, minimum height=7.5mm, fill=orange!10},
  arr/.style={-{Stealth}, thick}
]
  \node[box] (mask) {Mask $M$\\and pixels $\mathbf{I}(M)$};
  \node[box, right=7mm of mask] (mink) {Minkowski\\norm, order $p{=}6$};
  \node[box, right=7mm of mink] (feat) {Feature $\mathbf{f}\in\mathbb{R}^{8}$\\$\chi$, $o$, $\eta$};
  \node[box, fill=blue!8, below=8mm of feat] (std) {Standardize\\with training $\mu,\sigma$};
  \node[box, fill=blue!8, left=7mm of std] (lin) {Linear scores\\$W\tilde{\mathbf{f}}+b$};
  \node[box, fill=green!12, left=7mm of lin] (cls) {Class $\hat{k}$\\and margin};
  \draw[arr] (mask) -- (mink);
  \draw[arr] (mink) -- (feat);
  \draw[arr] (feat) -- (std);
  \draw[arr] (std) -- (lin);
  \draw[arr] (lin) -- (cls);
  \node[box, fill=black!5, below=7mm of mask] (pool) {Optional neck vector\\mask-pooled, 256-d};
  \node[box, fill=black!5, below=7mm of lin] (proj) {Projection to 32-d\\residual path};
  \draw[arr] (pool) -- (proj);
  \draw[arr] (proj) -- (lin);
\end{tikzpicture}
\caption{ICPH on one mask. The upper path is the chromatic feature that the estimate in Chapter~\ref{ch:system} actually trains. The lower path is the mask-pooled neck vector specified for the full segmenter. The linear layer is the prototype score. Both paths are cheap relative to the backbone.}
\label{fig:icph-block}
\end{figure}

\subsection{A per-mask illuminant that must not be trusted alone}
\label{sec:icph-illum}

Inside $M$ the diagonal model~\eqref{eq:diagonal} still applies, with one extra scalar. Shading and exposure multiply every channel by a positive $s(x)$, so
\begin{equation}
\mathbf{I}(x) = s(x)\, \mathbf{L} \odot \mathbf{R},
\qquad x\in M,
\label{eq:shaded}
\end{equation}
where the reflectance $\mathbf{R}$ has been pulled out of the pixel index because \Cref{ass:color} says the body is one color. Logos and drainage slots violate that pull-out on a small set of pixels. The head does not try to segment the logo. It relies on the Minkowski norm with $p>1$ to reduce the influence of the darkest slots, and on the training contamination described in Section~\ref{sec:icph-train} to see a strip of the neighboring crate.

The shades-of-gray estimate on the mask is~\eqref{eq:minkowski} with $\Omega=M$ and $p=6$:
\begin{equation}
\hat{\ell}_c
= \left(
\frac{1}{\lvert M\rvert}
\sum_{x\in M} \lvert I_c(x)\rvert^{6}
\right)^{1/6},
\qquad
\hat{\ell}_c \ge \varepsilon,
\label{eq:sog-mask}
\end{equation}
with $\varepsilon=10^{-6}$ only to avoid a division by zero on a degenerate mask. The corrected pixels are $\mathbf{I}'(x)=\mathbf{I}(x)\oslash\hat{\ell}$.

\begin{proposition}[A uniform face has no illuminant residual]
\label{prop:uniform}
Suppose $s(x)=s$ is constant on $M$ and $\mathbf{I}(x)=s\,\mathbf{L}\odot\mathbf{R}$ for every $x\in M$. Then $\hat{\ell}$ is parallel to $\mathbf{L}\odot\mathbf{R}$, the corrected color $\mathbf{I}'$ is constant across channels up to the common scale, and the chromaticity of $\mathbf{I}'$ is $(1/3,1/3,1/3)$ whenever no channel of $\mathbf{L}\odot\mathbf{R}$ is zero.
\end{proposition}

\begin{proof}
Under the hypothesis every pixel is the same vector $\mathbf{v}=s\,\mathbf{L}\odot\mathbf{R}$. The Minkowski norm of a constant set is the absolute value of that constant, so $\hat{\ell}=\lvert\mathbf{v}\rvert$ element-wise. Division by $\hat{\ell}$ yields a vector of ones on every channel that was nonzero. Its chromaticity, the normalization by the sum of channels, is the neutral point.
\end{proof}

\Cref{prop:uniform} is the reason ICPH does not classify $\mathbf{I}'$. On a clean crate face the classical correction succeeds at removing the illuminant and, in the same division, succeeds at removing the reflectance. Gray-world ($p=1$) and shades-of-gray ($p=6$) fail in the same direction. The failure is slightly worse for larger $p$ only because a larger $p$ concentrates on the brightest pixels, which on a uniform face are the same color as the mean. The estimate in Chapter~\ref{ch:system} is where this proposition becomes a table: both classical corrections fall to roughly one fifth accuracy, which is little above the one-in-six rate of a guess among six colors, as soon as the face is close to uniform.

The proposition also says where a residual can come from. If $s(x)$ varies, if a fraction of $M$ belongs to a neighbor, or if a logo has a different reflectance, then $\hat{\ell}$ is no longer parallel to a single $\mathbf{L}\odot\mathbf{R}$ and $\mathbf{I}'$ is no longer neutral. Those violations are exactly the nuisances of the passage. They are weak signals. A method that needs them in order to retain any color at all is a method that has thrown away $\mathbf{R}$. ICPH therefore keeps three views of the same pixels at once: the corrected chromaticity, which is informative when the mask is contaminated; the opponent coordinates of that chromaticity, which align with the amber--yellow and red--green confusions; and the chromaticity of $\hat{\ell}$ itself, which still carries $\mathbf{L}\odot\mathbf{R}$ and is the term the classical correction divides out.

\subsection{The eight-dimensional feature}
\label{sec:icph-feature}

Let $\bar{\mathbf{I}}'$ be the mean of the corrected pixels on $M$. Chromaticity drops the intensity:
\begin{equation}
\chi
= \frac{\bar{\mathbf{I}}'}{\mathbf{1}^{\top}\bar{\mathbf{I}}'+\varepsilon}
\in\mathbb{R}^{3}.
\label{eq:chroma}
\end{equation}
The opponent pair used here is the simplest linear map that separates red from green and that pair from blue,
\begin{equation}
o_{1}=\chi_{r}-\chi_{g},
\qquad
o_{2}=\tfrac{1}{2}(\chi_{r}+\chi_{g})-\chi_{b}.
\label{eq:opponent}
\end{equation}
The illuminant feature is the chromaticity of the Minkowski estimate, not its magnitude. Magnitude is brightness, and brightness is the fork shadow:
\begin{equation}
\eta
= \frac{\hat{\ell}}{\mathbf{1}^{\top}\hat{\ell}+\varepsilon}
\in\mathbb{R}^{3}.
\label{eq:eta}
\end{equation}
The head concatenates them,
\begin{equation}
\mathbf{f}
= \begin{bmatrix} \chi \\ o \\ \eta \end{bmatrix}
\in\mathbb{R}^{8}.
\label{eq:feature}
\end{equation}
\Cref{alg:feature} is the same construction in the order the implementation uses.

\begin{algorithm}[t]
\caption{Chromatic feature of one mask}
\label{alg:feature}
\begin{algorithmic}[1]
\Require pixels $\{\mathbf{I}(x):x\in M\}$, order $p=6$, floor $\varepsilon=10^{-6}$
\State $\hat{\ell}_c \gets \big(\lvert M\rvert^{-1}\sum_{x\in M}\lvert I_c(x)\rvert^{p}\big)^{1/p}$ for each channel $c$
\State $\hat{\ell} \gets \max(\hat{\ell},\varepsilon)$
\State $\mathbf{I}'(x) \gets \mathbf{I}(x)\oslash\hat{\ell}$
\State $\chi \gets \operatorname{mean}(\mathbf{I}') / (\mathbf{1}^{\top}\operatorname{mean}(\mathbf{I}')+\varepsilon)$
\State $o \gets (\chi_r-\chi_g,\ \tfrac{1}{2}(\chi_r+\chi_g)-\chi_b)$
\State $\eta \gets \hat{\ell}/(\mathbf{1}^{\top}\hat{\ell}+\varepsilon)$
\State \Return $\mathbf{f}\gets[\chi;\,o;\,\eta]$
\end{algorithmic}
\end{algorithm}

Two normalizations are doing different jobs and should not be collapsed. Division by $\hat{\ell}$ inside $\mathbf{I}'$ is a color-constancy step, and \Cref{prop:uniform} says it is destructive when it is used alone. Division by the channel sum inside $\chi$ and $\eta$ removes intensity after that step. The opponent map is not a normalization. It is a choice of axes. Amber and yellow differ mostly along a red--green direction that has been rotated slightly toward yellow; $o_1$ is a coarse version of that axis, not a perceptual uniform space. A more careful uniform color space would be justified if the prototypes were close. They are not close in the nominal palette. They become close only after a strong illuminant has moved the observations, which is the regime the linear classifier is trained for.

The nominal reflectances used to build the palette, written as linear RGB triples, are
\begin{equation}
\begin{aligned}
\mathbf{R}_{\text{amber}}&=(0.72,0.45,0.12),&
\mathbf{R}_{\text{green}}&=(0.15,0.55,0.22),&
\mathbf{R}_{\text{red}}&=(0.70,0.12,0.12),\\
\mathbf{R}_{\text{blue}}&=(0.12,0.25,0.62),&
\mathbf{R}_{\text{yellow}}&=(0.85,0.75,0.15),&
\mathbf{R}_{\text{white}}&=(0.82,0.82,0.80).
\end{aligned}
\label{eq:prototypes}
\end{equation}
These are design values, chosen to be separable in chromaticity and to resemble the families of returnable crates, not spectrophotometric measurements of a named manufacturer. Their chromaticities already predict the confusions. Amber and yellow both have a small blue coordinate and a large red coordinate. White is near $(1/3,1/3,1/3)$. Red and blue are far from everyone. The estimate later shows that pattern in the harsh-illuminant confusion matrix rather than as a surprise.

\begin{figure}[t]
\centering
\begin{tikzpicture}
\begin{axis}[
  width=0.78\textwidth,
  height=7.2cm,
  xlabel={$\chi_r$},
  ylabel={$\chi_g$},
  xmin=0, xmax=0.9,
  ymin=0, ymax=0.75,
  legend pos=north east,
  tick label style={font=\small},
  label style={font=\small},
  legend style={font=\small}
]
\addplot[only marks, mark=*, mark size=2.5pt] coordinates {
  (0.558,0.349)
  (0.163,0.598)
  (0.745,0.128)
  (0.121,0.253)
  (0.486,0.429)
  (0.336,0.336)
};
\node[font=\small, anchor=west] at (axis cs:0.558,0.349) {amber};
\node[font=\small, anchor=north west] at (axis cs:0.17,0.59) {green};
\node[font=\small, anchor=west] at (axis cs:0.745,0.128) {red};
\node[font=\small, anchor=west] at (axis cs:0.121,0.253) {blue};
\node[font=\small, anchor=south west] at (axis cs:0.486,0.429) {yellow};
\node[font=\small, anchor=east] at (axis cs:0.336,0.336) {white};
\end{axis}
\end{tikzpicture}
\caption{Chromaticity of the six design prototypes~\eqref{eq:prototypes}, plotted in the $(\chi_r,\chi_g)$ plane. Blue is omitted from the third coordinate and is small for amber and yellow, which is why those two classes approach each other when the illuminant estimate is noisy. White sits near the neutral point.}
\label{fig:chroma-plane}
\end{figure}

\subsection{Scores against prototypes}
\label{sec:icph-score}

A training set of features is standardized with the coordinate-wise mean $\mu\in\mathbb{R}^{8}$ and standard deviation $\sigma\in\mathbb{R}^{8}$ of that set,
\begin{equation}
\tilde{\mathbf{f}} = (\mathbf{f}-\mu)\oslash\sigma.
\label{eq:standardize}
\end{equation}
The score of class $k$ is linear,
\begin{equation}
z_k = W_{k,:}\,\tilde{\mathbf{f}} + b_k,
\qquad
W\in\mathbb{R}^{6\times 8}.
\label{eq:linear}
\end{equation}
The predicted class and the margin are
\begin{equation}
\hat{k}=\argmax_k z_k,
\qquad
\gamma = z_{\hat{k}} - \max_{k\neq\hat{k}} z_k.
\label{eq:margin}
\end{equation}
The margin is part of the product, not a training detail. A passage that contains one low-margin crate should be flagged for the operator even when the lattice is confident, because a swap of amber and yellow does not change the geometry that PLC checks. The system rule is given in Section~\ref{sec:sys-output}. The head itself only computes $\gamma$.

The training loss on a labeled feature is the cross-entropy of the softmax, plus an $\ell_2$ penalty on $W$,
\begin{equation}
\mathcal{L}
= -\log \softmax(z)_{y}
+ \frac{\lambda}{2}\lVert W\rVert_{F}^{2},
\label{eq:loss}
\end{equation}
with $\lambda=10^{-3}$ in the estimate. There is no penalty on $b$. The softmax is the usual one,
\begin{equation}
\softmax(z)_k
= \frac{\exp(z_k)}{\sum_{j\in\mathcal{K}}\exp(z_j)}.
\label{eq:softmax}
\end{equation}
This is a linear classifier on a hand-built feature. It is not a second detection network. The justification for stopping at a linear map is the dimension: eight coordinates, six classes, and a feature that was constructed so that the classes are roughly linear in it. If a field trial shows a curved boundary between amber and yellow, the replacement is a small multilayer perceptron on the same $\mathbf{f}$. The interface, the mask in and the class out, does not change.

The raw-RGB baseline that the estimate compares against ignores illumination. It takes the mean pixel $\bar{\mathbf{I}}$ on the mask, normalizes it to unit length, and chooses the prototype of nearest cosine,
\begin{equation}
\hat{k}_{\text{raw}}
= \argmax_{k\in\mathcal{K}}
\frac{\mathbf{R}_{k}^{\top}\bar{\mathbf{I}}}{\lVert\mathbf{R}_{k}\rVert\,\lVert\bar{\mathbf{I}}\rVert}.
\label{eq:raw}
\end{equation}
The gray-world and shades-of-gray baselines correct the pixels first, with $p=1$ and $p=6$, and then apply the same nearest-prototype rule to the corrected mean. Under \Cref{prop:uniform} those two baselines are predicted to collapse. The raw rule is predicted to be strong when $\mathbf{L}$ is near the white point and to decay as the random gain in the observation model grows. ICPH is predicted to decay more slowly, because $\eta$ still carries the direction of $\mathbf{L}\odot\mathbf{R}$ and the classifier can learn to read it. Chapter~\ref{ch:system} is where those predictions are checked. They are predictions about a model, and they are stated here so that the table cannot be mistaken for an exploration that found a pattern after the fact.

\subsection{The neck path that the full segmenter adds}
\label{sec:icph-neck}

The chromatic feature does not see the logo as a spatial pattern, and it does not see the difference between a crate face and a printed label of the same average color. The full head specified for YOLO-vf-seg therefore adds the mask-pooled neck vector. Let $U$ be the finest neck map that still has a channel count of $256$, and let $\Pi(M)$ be the set of spatial indices of $U$ that fall inside $M$. Mask pooling is
\begin{equation}
\mathbf{u}
= \frac{1}{\lvert\Pi(M)\rvert}
\sum_{q\in\Pi(M)} U_{q}
\in\mathbb{R}^{256}.
\label{eq:pool}
\end{equation}
A pointwise projection $P_{32}\in\mathbb{R}^{32\times 256}$ and a bias produce a 32-dimensional embedding. In parallel, a small perceptron maps the three raw channel means of the mask through a hidden width of 32 and out to the six classes, and that path is added to the linear scores of~\eqref{eq:linear}. The parameter count of this specification is
\begin{equation}
\begin{aligned}
N_{\text{proj}} &= 256\cdot 32 + 32,\\
N_{\text{mlp}} &= 3\cdot 32 + 32\cdot 32 + 32 + 32\cdot 6,\\
N_{\text{proto}} &= 32\cdot 6,\\
N_{\text{ICPH}} &= N_{\text{proj}}+N_{\text{mlp}}+N_{\text{proto}} = 9760,
\end{aligned}
\label{eq:params}
\end{equation}
where the prototype term counts a $32\times 6$ cosine table and the MLP term counts the illuminant residual path used in the cost model. Relative to a YOLOv8m-seg reference of $27.3\times 10^{6}$ parameters, the fraction is $3.6\times 10^{-4}$. The FLOPs of the head are one pooled vector per detection. They do not change the leading term of the segmenter, which for that reference model is on the order of $10^{2}$ GFLOPs. The estimate records the fraction and does not claim a new timing measurement on a device.

The split between what is trained in the estimate and what is specified for the system is deliberate. The Monte Carlo instrument trains~\eqref{eq:linear} on~\eqref{eq:feature}. It does not train a YOLO backbone, because no labeled passage images are part of this thesis. The neck path is specified so that a later training run on real frames has a place to put appearance that is not a global color, and so that the parameter budget is fixed before that run. Reporting the chromatic-path accuracy as if it were the full neck-path accuracy would be a misstatement. Chapter~\ref{ch:system} repeats the distinction next to the table.

\subsection{How the training distribution is built}
\label{sec:icph-train}

Each training face is drawn as follows. A class $k$ is uniform on the six prototypes. An illuminant gain is drawn as a log-normal vector,
\begin{equation}
\mathbf{g} = \exp(\boldsymbol{\xi}),
\qquad
\xi_c \sim \mathcal{N}(0,s^{2}),
\qquad
\mathbf{g}\gets \mathbf{g}/\operatorname{mean}(\mathbf{g}),
\label{eq:gain}
\end{equation}
with severity $s$ uniform between $0.05$ and $0.55$ during training. The normalization by the mean removes a pure exposure change, which shading already models, and keeps the draw on the illuminant direction. A shadow scalar $\rho$ is uniform on $[0.40,1]$, and a smooth shade multiplies the rows of the patch by a line from $0.72$ to $1.18$. With probability $0.25$ the first $18$ percent of the rows are replaced by a different prototype under the same $\mathbf{g}$ and $\rho$, which is the neighboring-crate strip. Pixel noise of standard deviation $0.012$ is added and negative values are clipped. The feature~\eqref{eq:feature} is computed on that patch, and the label is $k$.

Six thousand patches and five hundred full-batch gradient steps with step size $0.35$ are the training schedule of the estimate. The step size is large because the features are standardized and the model is linear; it is not a step size that would be copied onto a YOLO backbone. The $\ell_2$ penalty is applied to $W$ on every step. Nothing in this schedule is claimed to be optimal. It is fixed, it is seeded, and the test patches are drawn after training from the same generative steps at fixed severities that training never isolated.

The strip of a neighbor is the important augmentation. Without it the classifier can ignore $\chi$ and read the class from $\eta$ alone, because on a pure face $\eta$ is the chromaticity of $\mathbf{L}\odot\mathbf{R}$. With the strip, $\eta$ is pulled toward the neighbor and $\chi$ carries the disagreement between the two reflectances. That disagreement is what a real mask boundary looks like. Training without the strip would overstate accuracy on clean faces and understate the amber--yellow collision when the boundary is wrong.

\subsection{Calibration of the margin}
\label{sec:icph-margin}

The margin $\gamma$ in~\eqref{eq:margin} is not a probability. Softmax outputs of a linear model trained under a single generative severity are not calibrated to a different severity, and a warehouse shift is exactly a change of severity. The system therefore treats $\gamma$ as a rank, not as a chance. The operating rule, justified by the harsh-severity confusion in Chapter~\ref{ch:system}, is:

\begin{itemize}
\item if $\gamma$ is among the lowest tenth of margins seen on a calibration shift, the type is marked uncertain;
\item uncertain types are excluded from the automatic per-color tally and listed for the operator;
\item the total count from PLC is still reported, because geometry does not depend on $\gamma$.
\end{itemize}

A calibration shift is a few dozen pallets of known color at the start of a week, not a gray card. This is the one place the design asks a person to confirm a label. It is also the place where amber and yellow will concentrate if the lamps in that building are worse than the harsh setting of the estimate. The rule is a consequence of \Cref{prop:uniform} being only approximately false: when the residual that ICPH reads becomes small, the honest output is an abstention, not a forced name.

\subsection{Cost, interfaces, and the failures that remain}
\label{sec:icph-cost}

The runtime interface is a mask and a patch of pixels in, and a class plus a margin out. The training interface is a set of patches drawn as in Section~\ref{sec:icph-train}, or, when real images exist, the same feature computed on masks whose color a person has confirmed. The head does not modify the mask and does not modify the lattice. Those two statements are what keep the estimate separable. A color error cannot create a cell, and a gap-fill cannot change a prototype.

The failures that the derivation predicts, and that the estimate should be checked against, are the following.

\begin{enumerate}
\item At very small severity the raw cosine~\eqref{eq:raw} can beat ICPH, because the correction inside $\chi$ removes a reflectance signal that the raw matcher was using lawfully. A crossover at low severity is not a defect of the table. It is \Cref{prop:uniform} in mild form.
\item Amber and yellow remain the nearest pair. No linear feature built from a three-channel mean separates them once the gain noise exceeds their chromatic distance. More pixels do not help if the pixels are the same color. A logo-aware neck path might help. The chromatic path alone will not finish the job.
\item White is near neutral. Any error in $\hat{\ell}$ moves $\eta$ toward a chromatic class, and white is then confused with several classes rather than with one neighbor. The confusion matrix should show a white row that is spread out, not a white row that is swapped with a single partner.
\item A mask that contains more of the neighbor than of the crate will be classified as the neighbor. ICPH does not repair a bad mask. The segmenter remains responsible for $M$.
\end{enumerate}

These four points are the contract between this section and the color table. A result that contradicted them would mean the observation model and the head were not implementing the same mathematics. A result that matches them is evidence about the model, which is all a Monte Carlo estimate can be.

\subsection{Worked feature on two prototypes}
\label{sec:icph-worked}

The algebra is easier to see on a single pixel with no shade and no strip. Take a unit-mean illuminant $\mathbf{L}=(1.20,0.90,0.90)$ and the amber prototype. The observed color is
\begin{equation}
\mathbf{I}
= (1.20\cdot 0.72,\ 0.90\cdot 0.45,\ 0.90\cdot 0.12)
= (0.864,\ 0.405,\ 0.108).
\label{eq:worked-i}
\end{equation}
The Minkowski norm of a single pixel is the pixel, so $\hat{\ell}=\mathbf{I}$ and the corrected color is $(1,1,1)$. Chromaticity $\chi$ is neutral and both opponent coordinates are zero. The illuminant feature $\eta$ is the chromaticity of $\mathbf{I}$ itself,
\begin{equation}
\eta
= \frac{(0.864,0.405,0.108)}{1.377}
= (0.628,\ 0.294,\ 0.078).
\label{eq:worked-eta}
\end{equation}
The raw matcher sees $\mathbf{I}$, which is exactly $\mathbf{L}\odot\mathbf{R}_{\text{amber}}$, and still matches amber if no other prototype is closer after the gain. The classical corrected matcher sees white. ICPH sees a neutral $\chi$ and an $\eta$ that still points at amber under this lamp. That is the entire content of the head on a clean face: do not divide the reflectance away and then hope the classifier reconstructs it.

Now replace the blue channel's observation by a contamination, $\mathbf{I}=(0.864,0.405,0.250)$, as if a small amount of a brighter neighbor had entered the mean. Then $\hat{\ell}$ changes, $\chi$ leaves the neutral point, and $o_{2}$ becomes negative because the blue coordinate of $\chi$ rises. The linear layer has a coordinate, $o_2$, on which that contamination is visible. Whether the layer uses it is what the gradient steps decide. The hand calculation only shows that the coordinate exists, which \Cref{prop:uniform} says it does not exist on a pure face.

\subsection{What each block of the feature is worth}
\label{sec:icph-ablation}

\Cref{prop:uniform} says that $\chi$ and $o$ go to zero on a pure face, and that $\eta$ keeps $\mathbf{L}\odot\mathbf{R}$. The linear layer might still use the small residuals that shading and the neighbor strip leave in $\chi$. Whether it does is an empirical question about this feature, and it can be asked inside the same generator without a new theory. A second run, restarted from the same seed and the same schedule of six thousand patches and five hundred steps, trains six classifiers. Each classifier sees a subset of the eight coordinates. The subsets are the three chromaticity coordinates, the two opponent coordinates, the three illuminant coordinates, chromaticity with opponent, chromaticity with illuminant, and the full vector. Accuracy is measured at severities $0.05$, $0.30$, and $0.60$ on $2500$ fresh patches each. This run is not the run behind \Cref{tab:color}. Restarting the seed consumes the random stream in a different order, so the full-feature column here is a replicate, not a transcription.

\begin{table}[t]
\centering
\caption{Top-1 accuracy of linear heads trained on subsets of $\mathbf{f}$. The full row is a replicate of the ICPH feature, not the main table. Chance is $1/6\approx 0.167$.}
\label{tab:ablation}
\small
\begin{tabular}{@{}lccc@{}}
\toprule
Coordinates & $s=0.05$ & $s=0.30$ & $s=0.60$ \\
\midrule
Chromaticity $\chi$ only & $0.542$ & $0.534$ & $0.486$ \\
Opponent $o$ only & $0.538$ & $0.529$ & $0.480$ \\
Illuminant $\eta$ only & $0.864$ & $0.705$ & $0.504$ \\
$\chi$ and $o$ & $0.550$ & $0.532$ & $0.494$ \\
$\chi$ and $\eta$ & $0.976$ & $0.824$ & $0.592$ \\
Full $\mathbf{f}$ & $0.974$ & $0.827$ & $0.596$ \\
\bottomrule
\end{tabular}
\end{table}

\Cref{tab:ablation} is the practical form of the proposition. Chromaticity alone, the quantity a corrected matcher would trust, sits near $0.54$ at a mild lamp and falls toward $0.49$ at a harsh one. It is well above chance, because the shade ramp and the neighbor strip keep $\chi$ from being exactly neutral, and it is far below the raw matcher, which at low severity still sees $\mathbf{R}$ directly. Opponent coordinates alone do no better. They are a rotation of $\chi$, and a rotation cannot recover a signal the correction destroyed. Adding opponent to chromaticity moves the mild-lamp accuracy from $0.542$ to $0.550$. The opponent block is not where the accuracy lives.

The illuminant block alone is a different estimator. At $s=0.05$ it reaches $0.864$. On a nearly pure face $\eta$ is the chromaticity of $\mathbf{L}\odot\mathbf{R}$, and when $\mathbf{L}$ is near white that chromaticity is nearly the chromaticity of $\mathbf{R}$. At $s=0.60$ the same block falls to $0.504$, close to chromaticity-only. A strong gain rotates $\mathbf{L}\odot\mathbf{R}$ off the prototype, and $\eta$ rotates with it. The illuminant feature is a good feature exactly while the lamp is the nuisance the raw matcher can also tolerate, and it ceases to be sufficient at the severity ICPH was introduced to handle.

The combination of $\chi$ and $\eta$ is the head. At $s=0.05$ it reaches $0.976$, at $s=0.30$ it reaches $0.824$, and at $s=0.60$ it reaches $0.592$. The full eight-dimensional vector changes those three numbers to $0.974$, $0.827$, and $0.596$. Opponent coordinates are worth a few thousandths at the harsh end and a slight loss at the mild end, which is what an almost-redundant linear feature does under an $\ell_2$ penalty. They are kept in the specified head because they cost two coordinates and because the amber--yellow axis is an opponent axis. They are not kept because the ablation shows a large effect. A deployment that needs to shrink the feature can drop $o$ and expect to lose the margin visible in the last column, not the gap between ICPH and raw matching.

The ablation also warns against a story in which ``illumination invariance'' means the corrected chromaticity. The invariant-looking coordinates are the weak ones. The coordinates that carry the class are the ones the classical algorithm divides out, read by a classifier that has seen many lamps, together with the residual that contamination leaves in $\chi$. Invariance here is a property of the decision, not of a single normalized image. That is a weaker slogan and a more accurate description of \Cref{tab:ablation}.

\subsection{Severity as an angle, and the gradient the head follows}
\label{sec:icph-angle}

The severity $s$ is the standard deviation of a normal variable in log-gain space. It is worth converting to an angle so that ``$0.60$'' is not an arbitrary dial. Draw $\xi_c\sim\mathcal{N}(0,s^{2})$ independently and set $g_c=\exp(\xi_c)$, then renormalize $\mathbf{g}$ to unit sum or unit mean. The direction of $\mathbf{g}$ is the illuminant direction up to that normalization. For small $s$ the coordinates of $\xi$ are small and $\exp(\xi_c)\approx 1+\xi_c$. The component of $\xi$ orthogonal to $(1,1,1)$ has expected squared length $2s^{2}$, because the three coordinates have total variance $3s^{2}$ and the projection onto the white axis has variance $s^{2}$. The white vector has length $\sqrt{3}$, so the angle $\theta$ from white satisfies
\begin{equation}
\mathbb{E}[\theta^{2}]
\approx
\frac{2s^{2}}{3}.
\label{eq:angle}
\end{equation}
The root mean square angle is therefore $s\sqrt{2/3}$. At $s=0.05$ that is about $0.041$ radians, a little over two degrees. At $s=0.30$ it is about $0.24$ radians, near fourteen degrees. At $s=0.60$ the linearization is no longer quantitative: a one-standard-deviation log-gain of $0.60$ is a factor of about $1.8$ between channels, and the draws in the harsh column are lamp changes rather than tints. The mild column is a residual a locked camera might still see from shading. The harsh column is the size of change that separates daylight from a warm warehouse lamp.

The training update is the ordinary linear-softmax update, written out so that the ablation has a mechanism. With $p=\softmax(z)$ and label $y$, the gradient of the cross-entropy in~\eqref{eq:loss} with respect to the score is $p-e_{y}$. The gradient with respect to the row $W_{k,:}$ is
\begin{equation}
\frac{\partial\mathcal{L}}{\partial W_{k,:}}
= (p_k-\mathbf{1}_{k=y})\,\tilde{\mathbf{f}}^{\top}
+ \lambda W_{k,:}.
\label{eq:grad}
\end{equation}
A feature coordinate that is almost constant on the training set, as $\chi$ is on pure faces, produces a small outer product and a weight that the penalty $\lambda$ shrinks toward zero. A feature coordinate that tracks the class, as $\eta$ does when the lamp is mild and the face is pure, produces a weight that survives. The ablation is what~\eqref{eq:grad} does when some coordinates of $\tilde{\mathbf{f}}$ are noise of a destroyed reflectance and others are the reflectance the destruction was supposed to remove. No part of this argument requires the neck path. The neck path is an additional $256$-dimensional $\mathbf{u}$ that would enter~\eqref{eq:grad} through its own rows of $W$, and it is not part of the run.

\subsection{A numerical margin on the amber--yellow pair}
\label{sec:amber-number}

The confusion matrix is a count of swaps. The geometry of the swap is visible from the prototypes alone, before any classifier is trained. Take the amber and yellow triples in~\eqref{eq:prototypes}. Their cosine is
\begin{equation}
\frac{\mathbf{R}_{\mathrm{amber}}^{\top}\mathbf{R}_{\mathrm{yellow}}}
{\lVert\mathbf{R}_{\mathrm{amber}}\rVert\,\lVert\mathbf{R}_{\mathrm{yellow}}\rVert}
= \frac{0.9675}{0.8575\times 1.1435}
= 0.987.
\label{eq:cosine-ay}
\end{equation}
The numerator is $0.72\cdot 0.85+0.45\cdot 0.75+0.12\cdot 0.15$. The angle between the prototypes is $\arccos(0.987)\approx 9.3^{\circ}$. The same calculation for red against blue gives a cosine of $0.385$ and an angle near $67^{\circ}$. A gain perturbation of about fourteen degrees, the root-mean-square scale of $s=0.30$ under~\eqref{eq:angle}, is already larger than the amber--yellow separation and still smaller than the red--blue separation. The classifier can still do better than the cosine of the raw prototypes, because it sees $\eta$ and $\chi$ jointly and because the perturbation is random rather than adversarial, but it cannot manufacture distance that the reflectances do not have. Equation~\eqref{eq:cosine-ay} is why the abstention rule is aimed at a pair rather than applied uniformly, and why a deployment with a different palette should recompute this angle before trusting the harsh column of the table. A palette whose nearest pair sits at $30^{\circ}$ would make the same head look more accurate without any change in the algorithm. The accuracy in \Cref{tab:color} is an accuracy on~\eqref{eq:prototypes}, not an accuracy on color naming in the abstract.


\section{Perspective lattice completion}
\label{sec:plc}

\subsection{An error budget in crates, before the tables}

The tables measure a net error. The net is easier to read if it is split, even approximately, into the pieces the design can and cannot move. Consider a front face whose true occupancy is a solid $m_{\star}$ cells, depth $d$, truth $N=d\,m_{\star}$. Let the detector return $m_{\star}-q+f$ centers that snap into distinct true or false cells, where $q$ is the number of missed true cells and $f$ is the number of false cells that land outside the true set. The gap rule then adds $r$ restorations of true misses and $p$ phantom cells. The declared front count is
\begin{equation}
m = (m_{\star}-q+f) + r + p,
\label{eq:front-balance}
\end{equation}
and the booked total, when the depth reading is correct, is $\hat N = d\, m$. The signed front error is
\begin{equation}
m-m_{\star} = -q + f + r + p = -(q-r) + (f+p).
\label{eq:front-signed}
\end{equation}
The first parenthesis is unrestored misses. The second is false cells plus phantoms. PLC exists to make $r$ close to $q$ without making $p$ large. The span, by comparison, sets the declared count to the area of the bounding rectangle, which sets $r=q$ on a solid rectangle whose corners were detected, and sets $p$ equal to every truly empty cell inside that rectangle plus any cell the false positives pushed into the percentile span. On the stepped stack of \Cref{fig:grid}, $q$ may be zero and $p$ is already one for PLC and two for the span, before any detection noise. That is a structural $p$, not a random $p$.

When the depth reading is off by one, an additional $\pm m$ appears in $\hat N-N$. Section~\ref{sec:depth-sensitivity} estimates that term at about one crate of absolute error under the ten percent error rate, for fronts near ten cells. It does not enter~\eqref{eq:front-signed}. A table that shows PLC-single at mean absolute error $3.01$ is therefore a sum of at least three positive contributions: unrestored misses scaled by depth, phantoms and false cells scaled by depth, and depth-reading mistakes. Removing any one of them leaves the other two. This is why a three-tenth-of-a-crate improvement over the span can coexist with a three-crate absolute error, and why the improvement is still the right comparison. The span pays the same depth-reading term. The difference between the rows is the difference in $(q-r)$ and $(f+p)$, not the difference in the sensor.

Equation~\eqref{eq:front-signed} also fixes the sign of the bias. A rule with $r$ near $q$ and $p>0$ on stepped outlines has positive front bias. A rule with $r=0$, which is the raw mask count, has negative front bias equal to $-q+f$, and $q$ exceeds $f$ because the detection probabilities in~\eqref{eq:detprob} are all below one and the false-center mean is only $0.35$. Multiplying by $d\ge 1$ preserves the sign. The sign flip between the mask rows and the lattice rows of the counting table is~\eqref{eq:front-signed} with two different $(r,p)$, written before the table so that the flip is a prediction. Chapter~\ref{ch:system} checks it.

\subsection{The worked stack in the same letters}
\label{sec:letters}

The letters of~\eqref{eq:front-signed} are easier to trust if they are filled in on the one stack that was also counted by hand. The front truth has $m_{\star}=10$. The first frame produced nine detections, so if those nine landed in nine distinct true cells, then $q=1$ and $f=0$ before the gap rule. The rule restored the miss, so $r=1$, and it added the phantom beside the stepped top row, so $p=1$. Then
\begin{equation}
m-m_{\star} = -(1-1) + (0+1) = +1,
\label{eq:letters-num}
\end{equation}
and times a correct depth of $2$ the booked error is $+2$, which is the instrument's $22$ against $20$. The span on the same frame fills both empty cells of the top row, so its $p$ is $2$ if $q$ and $r$ cancel, and the front error is $+2$, times depth $24$ against $20$. The arithmetic does not use the random seed. It uses the grid and the visit order. A different visit order, as Section~\ref{sec:grid-drawn} noted, would have set $p=2$ for PLC as well and the two counters would have agreed at $24$. The estimated gap between them is partly an agreement about holes and partly an accident of left-to-right visitation on this particular step. On a solid rectangle with no step, structural $p$ is zero for both, and the gap that remains is only the random holes one rule restores and the other fills blindly.

The same letters explain a passage the operator should not treat as a success merely because the total moved closer to the truth. Suppose the detector missed two cells, $q=2$, the rule restored one, $r=1$, and invented one phantom, $p=1$, with $f=0$. The front error is $-(2-1)+(0+1)=0$. The booked total matches $d\,m_{\star}$ and the exact-match column of the table increments. The load is still wrong in two places: one crate was never seen and was not inferred, and one crate was inferred where the pallet is empty. The completed-cell count is $r+p=2$, which is the quantity that keeps this zero from looking like a clean observation. An exact total with a completed-cell count of zero is the passage the lattice actually saw. An exact total with a positive completed-cell count is a cancellation inside~\eqref{eq:front-signed}, and the record is required to show the cancellation rather than only the zero.

This is also the sense in which the exact-match rate of $0.337$ for PLC-single should not be read as ``one passage in three is correct for the right reason.'' Some of those exact passages are cancellations. The bias of $+0.83$ says the cancellations do not dominate: phantoms and false cells outweigh unrestored misses, but not by a full crate on average. A rate of exact matches near one third, a bias under one crate, and a mean absolute error of three crates are three different statements, and a shift supervisor who is given only the exact-match rate will think the system is three times more reliable than the mean absolute error says it is. The table reports all three because~\eqref{eq:front-signed} says they can move independently.


\subsection{The count is a property of the pallet, not of the mask set}
\label{sec:plc-count}

Let the front face be a set of cells $(r,c)$ with $r$ counted upward from the pallet and $c$ counted across the face. A cell is truly occupied or not. Write $O_{r,c}\in\{0,1\}$ and let $d\in\{1,2,3\}$ be the depth of the load in crates. Under the regularity that a missing front cell means the ranks behind it are missing as well,
\begin{equation}
N = d \sum_{r,c} O_{r,c}.
\label{eq:true-count}
\end{equation}
The image observes neither $d$ nor every $O_{r,c}$. It observes a noisy set of centers on the faces that the segmenter found, and the system observes a reading $\hat{d}$ of the depth. PLC estimates the occupied set and plugs it into~\eqref{eq:true-count}:
\begin{equation}
\hat{N} = \hat{d}\, \lvert\mathcal{F}\rvert,
\label{eq:plc-count}
\end{equation}
where $\mathcal{F}$ is the set of cells declared occupied. Everything else in this section is a specification of $\mathcal{F}$.

Equation~\eqref{eq:true-count} is where PLC differs from an area ratio. The area-ratio counter estimates each $O_{r,c}$ by a fraction of visible area and sums the fractions. PLC estimates each $O_{r,c}$ by a hard membership in $\mathcal{F}$. The hard decision is justified only because the pitch is known. If the pitch were unknown, the area ratio would be the more honest estimator, and \Cref{tab:gap} would be read the other way.

\begin{figure}[t]
\centering
\begin{tikzpicture}[
  font=\small,
  box/.style={draw, rounded corners=2pt, align=center, inner sep=4pt, minimum height=8mm},
  arr/.style={-{Stealth}, thick}
]
  \node[box] (c) {Mask centers\\in the image};
  \node[box, right=6mm of c] (h) {Homography\\to metres};
  \node[box, right=6mm of h] (o) {Median\\origin};
  \node[box, right=6mm of o] (s) {Snap to\\cells};
  \node[box, below=7mm of s] (g) {Fill gaps with\\two neighbors};
  \node[box, left=6mm of g] (d) {Multiply by\\depth $\hat{d}$};
  \node[box, fill=green!12, left=6mm of d] (n) {Count $\hat{N}$};
  \draw[arr] (c) -- (h);
  \draw[arr] (h) -- (o);
  \draw[arr] (o) -- (s);
  \draw[arr] (s) -- (g);
  \draw[arr] (g) -- (d);
  \draw[arr] (d) -- (n);
\end{tikzpicture}
\caption{PLC on one frame. Rectification, origin, and snapping use the known pitch. Gap completion is a single morphological step on the integer grid. Depth is an external reading. No step is learned.}
\label{fig:plc-block}
\end{figure}

\subsection{Rectification onto the front plane}
\label{sec:plc-homography}

A point $\mathbf{u}=(u,v,1)^{\top}$ in the image is sent to the front plane by the homography $H\in\mathbb{R}^{3\times 3}$,
\begin{equation}
\begin{bmatrix} X \\ Y \\ S \end{bmatrix}
= H \mathbf{u},
\qquad
\mathbf{p} = (X/S,\ Y/S).
\label{eq:homography}
\end{equation}
The two coordinates of $\mathbf{p}$ are metres on the face, with $Y$ upward and $X$ across, once $H$ has been scaled to the calibration target. $H$ is estimated at installation from a planar target of known size held where the front face will pass, using the direct linear transform and a robust rejection of mismatched corners~\cite{hartley2004mvg,fischler1981ransac,zhang2000calibration}. It is re-estimated when the camera is knocked. It is not re-estimated per truck. A truck that yaws enough to break planarity will show up as a rise in the residual of the snap, and that residual is a diagnostic the operator can see.

PLC never uses $H$ on pixels. It uses $H$ on centers. A center is the first moment of the mask, or the box center if the mask is empty. Using the center rather than the contour is a choice against the area ratio: the center of a half-visible crate still falls inside the right cell if the visible part is the inner part of the face, and it falls outside if the visible part is a thin sliver. The observation model of Chapter~\ref{ch:system} places isotropic Gaussian noise on the center, with standard deviations $3.5\,\mathrm{cm}$ across and $2.8\,\mathrm{cm}$ vertically, which is small relative to a $40\,\mathrm{cm}$ by $28\,\mathrm{cm}$ cell and large enough to move a center that was already on the boundary.

\subsection{Origin and integer cells}
\label{sec:plc-origin}

The pitch $(w,h)=(0.40,0.28)$ is known and the origin is not. The pallet is not fixtured to centimetre accuracy relative to the gate, and the truck does not stop at a mark. The origin is estimated from the detections themselves. For a set of rectified centers $\{(x_i,y_i)\}$,
\begin{equation}
\begin{aligned}
o_x &= \median_i \Big(
x_i - \big(\round(x_i/w - \tfrac{1}{2})+\tfrac{1}{2}\big)\, w
\Big),\\
o_y &= \median_i \Big(
y_i - \big(\round(y_i/h - \tfrac{1}{2})+\tfrac{1}{2}\big)\, h
\Big).
\end{aligned}
\label{eq:origin}
\end{equation}
Each term is the residual of a center after it has been snapped to the nearest cell of a lattice that was temporarily assumed to have origin zero. The median of those residuals is the shift that makes the lattice agree with the majority of the centers. The median, rather than the mean, is the robust choice: a false detection in the aisle contributes one residual and does not drag the origin, up to the breakdown point of a one-dimensional median, which is one half of the centers. A passage whose detections are more than half false is a failed segmentation, and no origin estimator repairs it.

Given $\mathbf{o}=(o_x,o_y)$, the integer cell of a center is
\begin{equation}
c = \round\!\left(\frac{x-o_x}{w}-\frac{1}{2}\right),
\qquad
r = \round\!\left(\frac{y-o_y}{h}-\frac{1}{2}\right).
\label{eq:snap}
\end{equation}
The half-cell offset matches the way the observation model generates centers, as the center of cell $(r,c)$ at $((c+\tfrac{1}{2})w,\ (r+\tfrac{1}{2})h)$ plus noise. Several centers may fall in one cell. PLC keeps the cell once. It does not weight the cell by the number of centers. A double detection of the same crate is therefore one crate, which is the behavior a counter wants, and a detection that straddles a boundary counts for whichever cell the rounded coordinate selects.

\begin{proposition}[Snap radius]
\label{prop:snap}
If the origin is exact and a center is perturbed by less than half a cell in each axis, \eqref{eq:snap} returns the true cell.
\end{proposition}

\begin{proof}
The unperturbed center of cell $(r,c)$ maps to the integers $(r,c)$ by~\eqref{eq:snap}. A perturbation smaller than $w/2$ in $x$ changes $(x-o_x)/w$ by less than $1/2$, so the rounded value does not change, and the same argument applies to $y$ with $h/2$.
\end{proof}

The hypothesis of \Cref{prop:snap} is the quantitative form of \Cref{ass:lattice}. The Gaussian noise used in the estimate has standard deviation $3.5\,\mathrm{cm}$ against a half-cell of $20\,\mathrm{cm}$ horizontally, so a miss of the snap from noise alone is rare. The errors that move the count are missed detections and false detections, not the rounding.

\subsection{One pass of gap completion}
\label{sec:plc-gap}

Let $\mathcal{D}$ be the set of cells that received at least one snapped center. Let $[r_0,r_1]\times[c_0,c_1]$ be the smallest integer rectangle containing $\mathcal{D}$, expanded by one cell in each direction and then clipped back to the cells that are within one step of that rectangle. The expansion gives the completion rule a place to stand; it does not add a row of phantom crates by itself, because a cell is added only by the neighbor test. Initialize $\mathcal{F}\gets\mathcal{D}$. For each empty cell $(r,c)$ in the rectangle, in row-major order,
\begin{equation}
\begin{aligned}
&n(r,c) = \mathbf{1}_{\{(r-1,c)\in\mathcal{F}\}}
+\mathbf{1}_{\{(r+1,c)\in\mathcal{F}\}}
+\mathbf{1}_{\{(r,c-1)\in\mathcal{F}\}}
+\mathbf{1}_{\{(r,c+1)\in\mathcal{F}\}},\\
&\text{if } n(r,c)\ge 2 \text{ then add } (r,c) \text{ to } \mathcal{F}.
\end{aligned}
\label{eq:gap}
\end{equation}
Diagonal neighbors do not count. A single pass is used, and the update is immediate, so a cell filled early in the pass can support a cell visited later. \Cref{alg:plc} writes the whole counter, including the depth multiplication.

\begin{algorithm}[t]
\caption{Perspective lattice completion on one frame}
\label{alg:plc}
\begin{algorithmic}[1]
\Require rectified centers $\{(x_i,y_i)\}$, pitch $(w,h)$, depth reading $\hat{d}$
\If{there is no center}
  \State \Return $0$
\EndIf
\State compute $(o_x,o_y)$ by~\eqref{eq:origin}
\State $\mathcal{D}\gets\emptyset$
\For{each center}
  \State compute $(r,c)$ by~\eqref{eq:snap} and insert $(r,c)$ into $\mathcal{D}$
\EndFor
\State $\mathcal{F}\gets\mathcal{D}$
\State $(r_0,r_1,c_0,c_1)\gets$ bounding rectangle of $\mathcal{D}$
\For{$r=r_0$ to $r_1$}
  \For{$c=c_0$ to $c_1$}
    \If{$(r,c)\notin\mathcal{F}$ and $n(r,c)\ge 2$}
      \State $\mathcal{F}\gets\mathcal{F}\cup\{(r,c)\}$
    \EndIf
  \EndFor
\EndFor
\State \Return $\hat{d}\cdot\lvert\mathcal{F}\rvert$
\end{algorithmic}
\end{algorithm}

The threshold of two orthogonal neighbors is the smallest number that distinguishes a hole in a wall from the end of a wall. One neighbor is an edge: every occupied cell on the boundary of a solid stack has empty cells beside it with exactly one occupied neighbor, and filling those would grow the stack by a rind. Two neighbors is a corner of a hole. Three or four would refuse to fill a hole that sits against the top of the load, which is where clipping happens. Two is the compromise, and it has a known over-fill: a stack whose true outline already has a step, such as a top row with only the right-hand crate present, presents the empty cell beside that crate with an upper or lower neighbor and a lateral neighbor. The rule fills the step. Section~\ref{sec:plc-example} computes that case on the grid used in the estimate.

Immediate update in row-major order makes the rule slightly path-dependent. A hole that is a whole missing row is not filled by a vertical pair of neighbors, because the cells of the missing row are visited when their vertical neighbors may exist but their horizontal neighbors do not yet, and a missing row has no horizontal support inside itself. A checkerboard cannot be completed into a solid block in one pass with a threshold of two if the occupied cells are only diagonal. Real pallet stacks are solid rectangles with a few missing cells, not checkerboards. The rule is designed for that shape. A checkerboard load would be a different product, and \Cref{ass:lattice} would have to be rewritten.

\subsection{Depth}
\label{sec:plc-depth}

The reading $\hat{d}$ is not inferred from the side image. The side image has already been used. Inferring depth from it would mean claiming that a front face determines the rank behind it, which is the mistake~\eqref{eq:true-count} is written to avoid. The estimate models the reading as
\begin{equation}
\hat{d} =
\begin{cases}
d & \text{with probability } 0.90,\\
\operatorname{clip}(d+\delta,1,3) & \text{with probability } 0.10,
\end{cases}
\qquad
\delta\in\{-1,+1\}
\text{ with equal probability.}
\label{eq:depth}
\end{equation}
Clipping at $1$ and $3$ matches the design range. A true depth of $1$ that is perturbed downward stays $1$. The probability $0.90$ is an assumption of the estimate, not a measured sensor specification. Changing it moves the bias of every method that multiplies by $\hat{d}$, which is all of them except the pure mask tally. The comparison among those methods remains informative because they share the reading.

When $\hat{d}$ is wrong, $\hat{N}$ is wrong by about $\lvert\mathcal{F}\rvert$, a whole rank. No gap rule repairs a bad depth. The system therefore reports $\hat{d}$ beside $\hat{N}$, so that an operator who knows the pallet was two deep can see a three. This is a display requirement in Section~\ref{sec:sys-output}, and it follows from the size of the error rather than from a preference about user interfaces.

\subsection{Why a union across the passage over-fills}
\label{sec:plc-passage}

A passage supplies $F$ frames, five in the estimate. The natural attempt to be robust is to snap every frame, count a vote $v(r,c)$ of how many frames produced cell $(r,c)$, and complete gaps on the cells whose vote is at least one. Call this PLC-passage. It is the temporal method that was estimated, and it is not the method the system should run.

The reason is the interaction of the vote with~\eqref{eq:gap}. A false center is rare in one frame. Across a union it is enough that each frame contribute a different false cell, or the same false cell just outside the true rectangle, for the bounding rectangle to grow and for the neighbor test to acquire support it did not have. Once a phantom cell is added, immediate update lets it support a second phantom. The bias has a sign: passage fusion adds cells and almost never removes them, because a vote threshold of one cannot delete a cell that a single noisy frame invented. The estimate measures that positive bias. The recommended counter is PLC on the single sharpest frame of the passage, PLC-single, where sharpness is the frame with the highest median margin on the color head and, if color was not run, the frame whose mask areas vary least across the stack.

A vote threshold of $\lceil F/2\rceil$ would block the union of one-frame false cells. It would also block a cell that was visible in only the sharpest frame and missed in the others, which is the usual signature of the bottom row. That trade is not estimated here. It is named so that the negative passage result is read as a result about the threshold of one, not as a result about the uselessness of a second frame. The second frame remains the right place to look when the first frame's residual of~\eqref{eq:origin} is large.

\subsection{Relation to the mask tally, the area ratio, and the span}
\label{sec:plc-baselines}

Four counters are the comparison set for PLC. All four, and both PLC variants, receive the same depth reading except the pure mask tally, which ignores depth and therefore estimates the front count rather than $N$.

The mask tally is the number of detected centers in the chosen frame,
\begin{equation}
\hat{N}_{\text{mask}} = \lvert\{\text{centers}\}\rvert.
\label{eq:mask-count}
\end{equation}
Multiplied by depth it becomes $\hat{N}_{\text{mask}\times d}=\hat{d}\,\lvert\{\text{centers}\}\rvert$. False positives and false negatives both enter with weight one. There is no notion of a cell.

The area ratio sums a noisy visible fraction for every crate the detector fires on, rounds, and multiplies by depth,
\begin{equation}
\hat{N}_{\text{area}}
= \hat{d}\;
\round\!\left(
\sum_{i\in\text{detected}} \alpha_i
\right),
\label{eq:area}
\end{equation}
with $\alpha_i$ drawn, in the estimate, from a normal distribution of mean $0.86$ and standard deviation $0.12$, clipped to $[0.35,1.15]$, plus a small false-positive area. The mean below one is the structural bias: partial masks contribute less than a crate and the counter never puts back a crate it did not detect.

The span counter uses the pitch and ignores occupancy inside the rectangle. With a robust width and height of the point cloud,
\begin{equation}
\begin{aligned}
\hat{n}_c &= \operatorname{clip}\!\left(
\round\!\big((\operatorname{pct}_{98}(x)-\operatorname{pct}_{2}(x))/w\big)+1,\,1,\,8\right),\\
\hat{n}_r &= \operatorname{clip}\!\left(
\round\!\big((\operatorname{pct}_{98}(y)-\operatorname{pct}_{2}(y))/h\big)+1,\,1,\,10\right),\\
\hat{N}_{\text{span}} &= \hat{d}\,\hat{n}_r\,\hat{n}_c.
\end{aligned}
\label{eq:span}
\end{equation}
The span is right on a solid rectangle whose extreme crates were detected, and it fills every hole, including holes that are truly empty. PLC is the span plus a refusal to fill cells that lack support, and plus a willingness to fill a single supported hole that the span also fills. On a nearly solid stack the two agree. They separate when the stack has a step, when false positives stretch the percentile, or when an interior miss is common enough that completing holes is better than trusting the raw set. The estimate is the measurement of that separation.

\subsection{Uncertainty}
\label{sec:plc-uncertainty}

PLC returns an integer. The uncertainty that should travel with it is not a softmax. Three quantities are available without a second model.

The origin residual, the median absolute residual inside~\eqref{eq:origin}, measures whether the centers agree on a lattice. A residual above a few centimetres means the snap is not in the regime of \Cref{prop:snap}.

The completion count $\lvert\mathcal{F}\setminus\mathcal{D}\rvert$ measures how much of the answer was inferred rather than seen. A passage that completes more than a small fraction of its cells is a passage whose bottom or top row was not really detected. The total can still be reported. It should be marked as completed, not as observed.

The depth bit, a Boolean that is false when the reading is known to be a default rather than a measurement, distinguishes a count that will be off by a rank if the default is wrong. The system displays all three next to $\hat{N}$. None of them is turned into a confidence interval in this thesis, because the observation model can compute an interval and a deployment cannot compute the observation model. The honest displayed uncertainty is the triple itself.

\subsection{A grid that over-fills by one cell}
\label{sec:plc-example}

The generative model used for the estimate can draw the front occupancy
\begin{equation}
O =
\begin{bmatrix}
1 & 1 & 1 \\
1 & 1 & 1 \\
1 & 1 & 1 \\
0 & 0 & 1
\end{bmatrix},
\label{eq:example-grid}
\end{equation}
with the first row at the bottom. Ten front cells are occupied. At depth $d=2$ the true count is $20$. The top row is a step: only the right-hand cell is present. Consider the empty top cell in the middle column, $(r,c)=(3,1)$, indexing from zero with row $3$ at the top. Its lower neighbor $(2,1)$ is occupied. Its right neighbor $(3,2)$ is occupied. The rule~\eqref{eq:gap} therefore adds $(3,1)$ even though $O_{3,1}=0$. The empty top-left cell $(3,0)$ has only its lower neighbor if it is visited before $(3,1)$ has been added, and in row-major order it is visited first, so it stays empty. The completion has invented one crate on the front face.

If, in the same frame, the detector missed one of the ten true cells and that cell has two occupied neighbors, the rule puts it back. The front count declared by PLC is then $10-1+1+1=11$, and times depth $2$ the estimate is $22$. The error is $+2$ on a truth of $20$. This is the worked passage reported numerically in Chapter~\ref{ch:system}: the same grid, a detected front set of nine cells in the first frame, PLC-single equal to $22$, and the span equal to $24$ because the span fills the bounding rectangle rather than one supported hole. The arithmetic is here so that the table is readable as a consequence of~\eqref{eq:gap} on~\eqref{eq:example-grid}, not as an isolated draw.

The same grid also shows why the bias of PLC-single can be positive even though the detector's bias is negative. The detector drops cells. The completion puts cells back, and on a stepped outline it puts back one cell too many. When the put-back exceeds the remaining misses, the signed error changes sign. A counter whose mean absolute error is better and whose bias has flipped from largely negative to slightly positive is a counter that has moved the error across zero, not a counter that has removed it. The operator-facing display of the completion count is what keeps that flip visible.

\subsection{Complexity}
\label{sec:plc-complexity}

Let $n$ be the number of centers in the frame and let $R\times C$ be the bounding rectangle, at most a few dozen cells for the loads in the design range. Origin estimation is a median of $n$ scalars, snapping is $O(n)$, and completion is $O(RC)$. There is no iteration in the sense of an optimizer and no parameter in the sense of a weight. The homography is a single matrix-vector product per center. Relative to the segmenter, PLC does not appear in a FLOPs budget. Relative to the area ratio it is the same order, a pass over the detections, and it uses a different summary of those detections.

The implementation consequence is that PLC can be re-run, inside the passage, on whichever frame the sharpness rule selects, without buffering features or resizing an image. The system in the next chapter uses that freedom. It keeps five masks only long enough to choose a frame, then runs \Cref{alg:plc} once.

\subsection{What PLC does not claim}
\label{sec:plc-claims}

PLC does not claim to see the rear rank. The factor $\hat{d}$ is an input. PLC does not claim to separate two crates that the segmenter merged into one center: a merged seam arrives as one point, snaps to one cell, and the neighboring cell is filled only if some other detection supports it. A systematic merge of every seam is not a hole pattern the rule recognizes. The vertical-feature gate~\eqref{eq:vf} remains the baseline's defense against that systematic merge. PLC does not claim invariance to the wrong pitch. A face of width $0.30\,\mathrm{m}$ snapped on a pitch of $0.40\,\mathrm{m}$ will alias, and \Cref{prop:snap} will not apply. A second family of crates requires a second pitch hypothesis and a selection between the two origins by the median residual. That selection is a small extension of~\eqref{eq:origin} and it is not in the estimate.

The claim that is made is narrower. On a load that satisfies \Cref{ass:lattice}, with centers that mostly fall in the correct cells and a depth reading of the quality in~\eqref{eq:depth}, declaring occupancy by a snapped set plus a two-neighbor completion, and multiplying by the reading, estimates $N$ with a smaller absolute error than counting masks, summing area ratios, or filling the bounding span. Chapter~\ref{ch:system} is the test of that claim under one observation model. The claim is not a field accuracy. The derivation in this section is what would have to be revisited if the test had come out the other way.

\subsection{The homography, written as a linear system}
\label{sec:dlt}

The map~\eqref{eq:homography} is used as a black box in the counter. The installation procedure has to compute it, so the linear system is part of the system design even though PLC never re-solves it per frame. A point $(X,Y)$ on the front plane and a pixel $(u,v)$ are related by
\begin{equation}
\begin{bmatrix} u \\ v \\ 1 \end{bmatrix}
\sim
H
\begin{bmatrix} X \\ Y \\ 1 \end{bmatrix},
\qquad
H=
\begin{bmatrix}
h_{11} & h_{12} & h_{13} \\
h_{21} & h_{22} & h_{23} \\
h_{31} & h_{32} & h_{33}
\end{bmatrix}.
\label{eq:H}
\end{equation}
The scale of $H$ is free. Clearing the homogeneous scale gives the two equations
\begin{equation}
\begin{aligned}
u(h_{31}X+h_{32}Y+h_{33}) &= h_{11}X+h_{12}Y+h_{13},\\
v(h_{31}X+h_{32}Y+h_{33}) &= h_{21}X+h_{22}Y+h_{23}.
\end{aligned}
\label{eq:dlt-pair}
\end{equation}
Stacked over $n\ge 4$ correspondences and moved to one side, the coefficients of the nine unknowns form a $2n\times 9$ matrix $A$ with a row pair
\begin{equation}
\begin{aligned}
&(X,\; Y,\; 1,\; 0,\; 0,\; 0,\; -uX,\; -uY,\; -u),\\
&(0,\; 0,\; 0,\; X,\; Y,\; 1,\; -vX,\; -vY,\; -v)
\end{aligned}
\label{eq:dlt-rows}
\end{equation}
for each correspondence. The solution is the singular vector of $A$ belonging to the smallest singular value~\cite{hartley2004mvg}. RANSAC wraps that solve: draw four correspondences, fit, count how many of the remaining points fall within a pixel threshold, and keep the set with the largest count~\cite{fischler1981ransac}. The target that supplies the correspondences is the planar board of the geometric calibration, whose spacing is known in metres, so the singular vector comes out in the scale PLC assumes when it divides by $w$ and $h$.

Four non-collinear correspondences determine $H$. The target provides many more so that a mismatched corner can be rejected. The condition the installer checks is not the algebraic residual of $A$. It is the geometric residual on the rectified plane, in centimetres, because \Cref{prop:snap} is stated in centimetres. A fit that is excellent in pixels and scaled wrong by ten percent will alias the lattice. The tape check on $(w,h)$ and the residual check on $H$ are the same check read from two sides.

The same linear system says what a yaw of the truck does. A modest yaw of a still-planar face is absorbed into $H$ if $H$ is re-estimated, and it is a bias on every center if $H$ is the one computed at installation for a face that was square to the camera. The design does not re-estimate $H$ per truck. It relies on the gate geometry to keep the yaw inside the snap radius. A truck that approaches at a large angle violates that reliance, and the origin residual grows. The residual is the detector of the violation. It is not a substitute for re-estimating $H$ on a truck that is known to arrive askew. If a building's traffic pattern makes large yaw routine, the installation procedure changes and the mathematics of~\eqref{eq:dlt-rows} does not.

\subsection{Where the detection probability spends the misses}
\label{sec:detprob-discuss}

The snap and the gap rule see only the centers that survive~\eqref{eq:detprob}. The shape of that probability is the reason a uniform hole-filler would be the wrong rule. With $\beta=0$ the bottom row is kept with probability $0.70$, the top row with probability $0.84$, and a middle row with probability $0.96$. With $\beta=1$ those numbers fall by $0.18$, to $0.52$, $0.66$, and $0.78$, before the clip at $0.35$ ever binds. The clip binds only if a future model adds a further penalty. In the stated model the bottom row at full blur is still detected more often than not.

\begin{figure}[t]
\centering
\begin{tikzpicture}
\begin{axis}[
  width=0.86\textwidth,
  height=6.6cm,
  xlabel={Blur variable $\beta$},
  ylabel={Detection probability},
  ymin=0.4, ymax=1.0,
  xmin=0, xmax=1,
  legend style={font=\small, at={(0.5,-0.24)}, anchor=north, legend columns=3},
  tick label style={font=\small},
  label style={font=\small},
  grid=major,
  grid style={gray!25}
]
\addplot+[mark=*] coordinates {(0,0.70) (0.5,0.61) (1,0.52)};
\addplot+[mark=square*] coordinates {(0,0.84) (0.5,0.75) (1,0.66)};
\addplot+[mark=triangle*] coordinates {(0,0.96) (0.5,0.87) (1,0.78)};
\legend{Bottom row, Top row, Middle row}
\end{axis}
\end{tikzpicture}
\caption{Detection probability~\eqref{eq:detprob} against blur, for a stack tall enough to have a distinct bottom, middle, and top. The gap rule is asked to repair the bottom curve, not the middle one.}
\label{fig:detprob}
\end{figure}

On a front face of four rows and three columns with every cell occupied, twelve crates, the expected number of misses at $\beta=0$ is
\begin{equation}
3(1-0.70)+3(1-0.84)+6(1-0.96)
= 0.90+0.48+0.24
= 1.62.
\label{eq:expected-miss}
\end{equation}
At $\beta=1$ the same sum is $3(0.48)+3(0.34)+6(0.22)=3.78$. A middle miss is the miss the two-neighbor rule is built to restore: the cell above and the cell below exist. A bottom miss has no lower neighbor. It is restored only when both lateral neighbors were detected, which on a three-column bottom row is true for the middle cell and only half-supported for the two ends. The rule is therefore a better repair of an interior hole than of a shadowed corner. Equation~\eqref{eq:expected-miss} says there is about one and a half misses to repair on a clean frame of that size, and nearly four on a blurred one. False positives arrive at a Poisson mean of $0.35$, an order of magnitude smaller than the blur-induced misses and large enough, across five independent frames, to supply the phantom cells that PLC-passage unions. The asymmetry is the design reason for trusting one frame.

The expected miss count is not the expected counting error. A miss that the gap rule restores costs nothing. A miss at a corner that the rule cannot see costs $\hat{d}$. A false positive that lands in an empty cell and acquires a second neighbor costs $\hat{d}$ the other way. The tables in Chapter~\ref{ch:system} are the net of those three effects. The expectation~\eqref{eq:expected-miss} only explains why the net is not zero and why it depends on which row was lost.

\subsection{The worked grid, drawn}
\label{sec:grid-drawn}

\Cref{fig:grid} draws~\eqref{eq:example-grid} before and after the pass~\eqref{eq:gap}. The occupied cells are the ten true front crates. The cell that the rule adds is the middle of the top row, which has a neighbor below and a neighbor to the right. The corner cell of that row has, at the moment it is visited, only the neighbor below, so it stays empty. The drawing is the whole of the positive bias on this stack: one invented front cell, times a depth of two, is two invented crates, before any missed detection is put back. When the first frame of the instrument also misses one true cell and the rule restores it, the invented cell remains and the declared front occupancy is eleven. Eleven times two is the $22$ reported later against a truth of $20$.

\begin{figure}[t]
\centering
\begin{tikzpicture}[
  font=\footnotesize,
  c/.style={draw, minimum width=8mm, minimum height=7mm, inner sep=0pt, anchor=center},
  on/.style={c, fill=orange!35},
  off/.style={c, fill=black!6},
  add/.style={c, fill=green!25}
]
  \node at (1.2,3.5) {(a) true front face};
  \foreach \r in {0,1,2} {
    \foreach \col in {0,1,2} {
      \node[on] at (\col*0.85, \r*0.75) {};
    }
  }
  \node[off] at (0, 2.25) {};
  \node[off] at (0.85, 2.25) {};
  \node[on] at (1.70, 2.25) {};
  \node[rotate=90, anchor=south] at (-0.55,1.1) {bottom $\rightarrow$ top};

  \begin{scope}[shift={(5.4,0)}]
    \node at (1.2,3.5) {(b) after one pass};
    \foreach \r in {0,1,2} {
      \foreach \col in {0,1,2} {
        \node[on] at (\col*0.85, \r*0.75) {};
      }
    }
    \node[off] at (0, 2.25) {};
    \node[add] at (0.85, 2.25) {};
    \node[on] at (1.70, 2.25) {};
    \node[align=center] at (1.2,-0.7) {shaded cell was empty\\and has two neighbors};
  \end{scope}
\end{tikzpicture}
\caption{The stepped stack of the worked example. The top row truly contains one crate. Gap completion fills the cell beside it and does not fill the corner, because the corner is visited first and has only one orthogonal neighbor at that moment.}
\label{fig:grid}
\end{figure}

The order dependence is visible in the drawing and is easy to overstate. If the pass visited columns from right to left, the middle top cell would still have its lower neighbor and its right neighbor, both original, and would still be filled. The corner would then be visited after the middle cell existed, and would gain a lateral neighbor. A right-to-left pass would fill both empty cells. The implementation visits left to right and therefore fills one. This is a reason to treat the threshold and the visit order as part of the algorithm that was estimated, not as a cosmetic of the pseudocode. A later change to the visit order is a different counter and needs its own row in the table. The present row is the left-to-right, immediate-update, threshold-two counter.

\subsection{A bound on a single false center}
\label{sec:false-center}

One false center can do only a limited amount of damage on a single frame, which is the other half of the argument against the passage union. Suppose the true occupied set is a solid rectangle and the origin is exact. A false center falls either in an occupied cell, in which case the snap ignores the duplicate, or in an empty cell. An empty cell adjacent to the rectangle has at most one orthogonal neighbor inside the rectangle if it shares a side, and it has none if it only touches a corner. The test $n\ge 2$ does not add a side-adjacent cell on the strength of one false center. The false center occupies a cell by itself, and the count grows by $\hat{d}$, once. It does not start a rind.

The passage union removes that limit. Frame $1$ may place a false center on the left of the rectangle and frame $2$ a false center above that false cell, or two frames may place false centers in the same empty side cell and in a second empty cell that becomes its neighbor through the immediate update. Vote threshold one treats each of those as support. The single-frame bound, ``one false center adds at most one cell,'' is false of the union. The estimated bias of $+5.33$ crates for PLC-passage against $+0.83$ for PLC-single is the bound failing in the direction the mechanism predicts. It is also the reason the software specification has no vote accumulator. An accumulator with threshold one is not a harmless extra. It is the counter the estimate rejected.

\subsection{Why the span is the serious baseline}
\label{sec:span-serious}

The mask tally and the area ratio are easy to beat once depth is multiplied in, and a comparison that stopped there would overstate the lattice. The span is the serious baseline because it already knows $(w,h)$ and it already refuses to count raw detections. The three tenths of a crate that separate PLC-single from the span are the entire geometric contribution of the neighbor test, and they should be described at that size.

On a solid rectangle whose corner crates were detected, the span's percentile width and height recover the true number of rows and columns, and the span books $d$ times that product. PLC books the same product if every interior hole was restored and no phantom was added. The two counters differ when the outline is not a solid rectangle, when a false center stretches a percentile, or when an interior miss is common and the span, which does not look inside the rectangle, cannot tell a hole from a filled cell. The stepped top row is the first case. A false center just outside the left column is the second: the second percentile of $x$ moves, $\hat n_c$ in~\eqref{eq:span} can increase by one, and the span multiplies a whole extra column by the number of rows and by $\hat d$. PLC snaps that false center into one cell. If the cell has fewer than two neighbors it stays a single phantom, not a column. The neighbor test is a limit on how much one false center can inflate the count, which is the same limit Section~\ref{sec:false-center} stated from the other direction, and it is the reason a method with a slightly higher exact-match rate can still be the one to deploy when the alternative's errors come in whole columns.

None of this makes the span a straw man. On the loads the generator actually draws, most faces are nearly solid, the span's mean absolute error is $3.31$ against PLC-single's $3.01$, and a deployer who implemented only the span would have captured most of the gain over mask-times-depth. The thesis still specifies the neighbor test, because the extra code is a double loop over a few dozen cells, because the bias drops from $+1.60$ to $+0.83$, and because the failure mode of an extra column is worse than the failure mode of a single phantom cell on a step. The specification would be different if the estimate had gone the other way. It did not.

\begin{figure}[p]
\centering
\begin{tikzpicture}[
  font=\small,
  box/.style={draw, rounded corners=2pt, align=center, inner sep=5pt, minimum width=32mm, minimum height=12mm},
  arr/.style={-{Stealth}, thick}
]
  \node[box, fill=black!6] (det) {Detections\\$m_{\star}-q+f$};
  \node[box, fill=green!15, right=14mm of det] (rule) {Gap rule\\$+r$ restorations\\$+p$ phantoms};
  \node[box, fill=orange!15, right=14mm of rule] (front) {Front count\\$m$};
  \node[box, fill=blue!10, below=16mm of front] (book) {Booked total\\$\hat d \cdot m$};
  \node[box, fill=red!8, left=14mm of book] (depth) {Depth reading\\correct, or $\pm 1$};
  \node[box, fill=black!4, left=14mm of depth] (truth) {Truth\\$N=d\,m_{\star}$};
  \draw[arr] (det) -- (rule);
  \draw[arr] (rule) -- (front);
  \draw[arr] (front) -- (book);
  \draw[arr] (depth) -- (book);
  \draw[arr, dashed] (truth) -- (book);
  \node[align=center, font=\footnotesize] at ($(det.south)+(0,-0.9)$) {misses $q$,\\false cells $f$};
  \node[align=center, font=\footnotesize] at ($(truth.south)+(0,-0.85)$) {signed error\\$\hat N-N$};
\end{tikzpicture}
\caption{How a booked count is assembled, in the letters of~\eqref{eq:front-signed}. The detector contributes misses and false cells. The gap rule puts some misses back and may add phantoms. Depth scales the front count and, when the reading is wrong, adds or removes a whole rank. The dashed arrow is the comparison the estimate computes. It is not an input to the system.}
\label{fig:balance}
\end{figure}


\subsection{How the two updates meet}
\label{sec:meet}

ICPH and PLC share the mask and nothing else. The mask provides pixels to the feature~\eqref{eq:feature} and a center to the snap~\eqref{eq:snap}. After that split, a color error and a count error are different events. The per-color count that \Cref{def:task} asks for is the combination
\begin{equation}
\hat{N}_{k}
= \hat{d}\;
\lvert\{ (r,c)\in\mathcal{F} : \hat{k}(r,c)=k \}\rvert,
\label{eq:per-color}
\end{equation}
where a completed cell that no detection owns inherits the color of the majority of its orthogonal neighbors, and a tie inherits the color of the neighbor with the larger margin. Inheritance is specified here and not estimated, because the Monte Carlo trials do not assign classes and cells in the same draw. The specification is the one a deployment would code. Its error, when a completed cell inherits a swapped amber label, is the sum of the two errors the estimate does measure. An upper bound that does not pretend to be tighter than the data is
\begin{equation}
\lvert\hat{N}_{k}-N_{k}\rvert
\le \lvert\hat{N}-N\rvert + E_{k},
\label{eq:bound}
\end{equation}
where $E_k$ is the number of front crates whose type was swapped into or out of class $k$, times $\hat{d}$ if the depth reading was applied to both. The bound says that type errors and count errors add. It does not say they cancel. The system reports them as separate lines so that a cancellation in the total cannot hide a swap in the colors.

The two updates are also independent in their requirements. ICPH needs \Cref{ass:color} and a mask that overlaps the face. PLC needs \Cref{ass:lattice}, a homography, and a depth reading. A building that controls its lamps and not its geometry would enable PLC and could turn ICPH down to the raw matcher. A building that has irregular mixed cartons and stable lamps would do the opposite, and would leave this thesis's lattice behind in favor of the area ratio. The brewery passage, as defined in Section~\ref{sec:scene}, enables both, which is why both are in the updated model.

The updated model, named in full, is YOLO-vf-seg with the vertical-feature neck~\eqref{eq:vf}, the prototype mask assembly~\eqref{eq:yolact}, the chromatic head of \Cref{alg:feature} and~\eqref{eq:linear}, and the lattice counter of \Cref{alg:plc}. The next chapter places that model in the passage: where the camera stands, what the computer stores, how a frame becomes a line on the operator's tally, and what the Monte Carlo estimate says about the head and the counter when the passage is simulated rather than filmed.

\chapter{System Construction and Estimation}
\label{ch:system}

The modules of Chapter~\ref{ch:methods} are not yet a system. A passage has a camera, a trigger, a computer, a depth reading, and a person who has to act on a wrong tally. Section~\ref{sec:build} specifies that system around YOLO-vf-seg, ICPH, and PLC, at the level of geometry, data flow, calibration, timing, and the operator's record. Section~\ref{sec:estimate} then estimates the two modules under a Monte Carlo model whose parameters are the same numbers used in the specification. The estimate is a property of the model. It is the evidence this thesis can offer without a filmed passage, and the text keeps that limit visible next to every table.

\section{System construction}
\label{sec:build}

\subsection{The gate}
\label{sec:gate}

The system is installed at a warehouse door that every finished pallet passes, rather than on the forklift. A camera on the truck would move with the load and would see the load from a fixed relative pose, which is attractive, and it would require power, vibration isolation, and a radio on every truck, which this design refuses. A gate camera sees many trucks. The price is that the face is in view for a short interval, and that the interval depends on the driver.

The nominal installation is a camera whose optical axis is horizontal, at a height of $1.4\,\mathrm{m}$ above the floor, aimed perpendicular to the path of the truck at a standoff of about $2.5\,\mathrm{m}$ from the plane in which the front face is expected. At that standoff a face of width $1.6\,\mathrm{m}$, four nominal crates, subtends a moderate field of view, and a stack of six crates at $0.28\,\mathrm{m}$ reaches $1.68\,\mathrm{m}$ above the pallet. A pallet height of roughly $0.15\,\mathrm{m}$ puts the top of a six-high stack near $1.8\,\mathrm{m}$. The camera height of $1.4\,\mathrm{m}$ looks slightly down on the lower rows and slightly up at the top row, which is the perspective the homography is asked to remove. A camera placed at the top of the stack would see the top faces and would lose the bottom row to the pallet. A camera placed at fork height would see the forks and would clip the top. The middle of the design range is the compromise.

The shutter is specified from the truck speed rather than from a frame-rate fashion. At $1.5\,\mathrm{m\,s^{-1}}$, a walking pace, a $2\,\mathrm{ms}$ exposure moves the load by $3\,\mathrm{mm}$ during the exposure. That smear is small beside the $3.5\,\mathrm{cm}$ center noise already admitted in the lattice. At $5\,\mathrm{m\,s^{-1}}$ the same exposure moves the load by $10\,\mathrm{mm}$, still inside the half-cell of \Cref{prop:snap}, and a $10\,\mathrm{ms}$ exposure would move it by $50\,\mathrm{mm}$, which is a quarter of a cell and no longer negligible. The specification is therefore an exposure at or below $2\,\mathrm{ms}$, with the gain raised to compensate, and a frame rate of $15\,\mathrm{Hz}$ so that a face which is fully in view for half a second contributes about seven frames, of which the estimate's five-frame passage is a conservative subset. The blur variable in the observation model is the normalized version of whatever this exposure fails to prevent. It is not a second shutter design.

Illumination is not redesigned. The gate inherits the dock. What the system specifies instead is a prohibition and a calibration. The prohibition is against an automatic white balance in the camera firmware. Firmware white balance is a global gain of the kind Section~\ref{sec:color-principles} rejected, and it moves from frame to frame as the truck's painted mast enters the field. The camera is locked to a fixed gain and a fixed color balance for the shift. ICPH then sees a stable $\mathbf{L}$ inside the passage and a varying $\mathbf{L}$ between shifts, which is the variation the head was built for. The calibration is the weekly set of known-color pallets used to rank the margin, as in Section~\ref{sec:icph-margin}.

\subsection{What is computed, and where}
\label{sec:compute}

One industrial PC at the gate is sufficient. The segmenter is the heavy step. ICPH and PLC are not. A medium YOLO-seg model on a current embedded GPU sustains a frame rate above the $15\,\mathrm{Hz}$ of the camera if the input is kept near the model's training resolution; the warehouse result cited in Section~\ref{sec:warehouse-yolo}, $84.8$ frames per second for a small variant, is evidence that the family has margin, not a measurement of this PC. The design therefore budgets the segmenter at one frame period, the color head at well under a millisecond for a few dozen masks, and the lattice at well under a millisecond. The budget that can actually break is the segmenter. If a deployment finds that the medium model misses the frame period, the specified downgrade is a smaller YOLO-vf-seg variant with the same two heads, not the removal of PLC. Removing PLC reintroduces the factor-of-depth error, which is larger than the accuracy difference between a small and a medium backbone.

\begin{figure}[t]
\centering
\begin{tikzpicture}[
  font=\footnotesize,
  box/.style={draw, rounded corners=2pt, align=center, inner sep=3.5pt, minimum height=7mm},
  arr/.style={-{Stealth}, semithick}
]
  \node[box, fill=black!8] (cam) {Locked\\camera};
  \node[box, right=5mm of cam] (trig) {Passage\\trigger};
  \node[box, right=5mm of trig] (buf) {Frame\\buffer, $F{\le}7$};
  \node[box, right=5mm of buf] (yolo) {YOLO-vf-seg\\masks, centers};
  \node[box, fill=orange!15, below=6mm of yolo] (icph) {ICPH\\class, margin};
  \node[box, fill=green!12, left=5mm of icph] (pick) {Pick the\\sharp frame};
  \node[box, fill=green!12, left=5mm of pick] (plc) {PLC\\on that frame};
  \node[box, left=5mm of plc] (dep) {Depth\\reading $\hat{d}$};
  \node[box, fill=blue!8, below=6mm of pick] (rec) {Tally record\\per color and total};
  \draw[arr] (cam) -- (trig);
  \draw[arr] (trig) -- (buf);
  \draw[arr] (buf) -- (yolo);
  \draw[arr] (yolo) -- (icph);
  \draw[arr] (icph) -- (pick);
  \draw[arr] (pick) -- (plc);
  \draw[arr] (dep) -- (plc);
  \draw[arr] (plc) -- (rec);
  \draw[arr] (icph) -- (rec);
\end{tikzpicture}
\caption{Data flow of one passage. The camera's white balance is locked. The buffer holds a handful of frames. The segmenter runs on each. ICPH scores every mask. PLC runs once, on the frame the margins select. The depth reading joins only at the multiplication.}
\label{fig:dataflow}
\end{figure}

The software is a single process with the stages of \Cref{fig:dataflow}. The trigger is a motion detector or a break-beam at the gate; the design does not depend on which, provided the beam fires before the face is centered and holds the buffer open until the face has left. Frames that fail a coarse brightness check are dropped rather than segmented. On each kept frame the segmenter writes masks, boxes, and centers. ICPH writes a class and a margin per mask. When the trigger closes, the frame with the highest median margin is passed to PLC together with the depth reading. The record in Section~\ref{sec:sys-output} is written once. There is no tracking identity across frames in the specified system. Section~\ref{sec:plc-passage} is why: a union of cells made the count worse under the model, and the software does not implement a method the estimate has already rejected.

The depth reading arrives on a separate path. The preferred source is the pallet build record, if the packaging line already knows whether it stacked one, two, or three deep. The fallback is a manual default for the shift, displayed as a default so that the uncertainty triple of Section~\ref{sec:plc-uncertainty} has its depth bit clear. The system does not infer depth from the side masks. That constraint is enforced in the software by the absence of any code path from the mask area to $\hat{d}$.

\subsection{The segmenter as a component}
\label{sec:segmenter-component}

YOLO-vf-seg is used as specified in Section~\ref{sec:problem-gap}. The input is a three-channel image at the network's native resolution. The neck applies the vertical-feature gate~\eqref{eq:vf} on the three pyramid levels that feed the head. The detection head emits boxes and objectness. The mask head emits thirty-two prototypes and a coefficient vector per detection, combined by~\eqref{eq:yolact} and cropped to the box. Non-maximum suppression keeps one detection per crate when the boxes overlap heavily. The suppression threshold is part of the deployment calibration, not part of the Monte Carlo model. The model represents a miss by a detection probability that already folds suppression, shadow, and blur into one number. A future measurement on real frames would have to unfold that number before anyone tuned the suppression threshold against this thesis's tables.

The segmenter is not trained in this thesis. A deployment train would start from a public YOLO-seg checkpoint trained on COCO~\cite{lin2014coco,jocher2023yolov8} and would continue on gate images labeled for masks and for the six colors. The color labels are the same labels ICPH needs, so one labeling pass serves both the mask head's class branch and the chromatic head. The vertical-feature gate can be trained from the start or folded in after the baseline converges. Neither schedule is estimated here. What the system specification fixes is the interface the training must preserve: masks whose first moment is a usable center, and pixels inside the mask that are still the crate's pixels. A training run that inflates masks to win a region overlap will hand ICPH the neighbor's color and will hand PLC a center that has drifted toward the neighbor. The overlap loss is therefore not the only training criterion a deployment should watch. The center error in centimetres on the rectified face, the quantity \Cref{prop:snap} bounds, is the criterion that matches the lattice.

\subsection{Calibration}
\label{sec:calibration}

Three calibrations are specified, and they are not optional in the way a software default is optional.

The geometric calibration estimates $K$ and $H$. A planar target of known spacing is held in the plane of a typical front face, the camera is locked, and $H$ is computed by the direct linear transform with RANSAC rejection of corners that the detector mismatched~\cite{zhang2000calibration,hartley2004mvg,fischler1981ransac}. The scale of $H$ is fixed by the target, in metres. After calibration, a person checks that the residual of a crate-sized rectangle on the target is small beside $w/2$. This calibration is repeated when the camera moves. A temperature drift of a fixed camera is ignored until that residual, checked at the start of a week, grows.

The pitch calibration is a confirmation, not an estimation. The nominal $(w,h)=(0.40,0.28)\,\mathrm{m}$ is checked with a tape on the crate family that the gate will see. If a second family is present, its pitch is entered as a second hypothesis. The origin estimator~\eqref{eq:origin} then runs twice and the hypothesis with the smaller median absolute residual is kept. The estimate in Section~\ref{sec:estimate} uses one pitch. The second hypothesis is specified so that the software has a place to put it, and so that nobody ``solves'' a second family by stretching $H$.

The color calibration is the weekly set of known pallets. It does not refit $W$. It records the distribution of margins and sets the abstention threshold at the tenth percentile of that distribution, as Section~\ref{sec:icph-margin} states. Refitting $W$ is a maintenance action, taken when the abstention rate rises for a week, not a per-passage action. Refitting on the passage itself would let a single badly lit truck move the prototypes.

\subsection{Timing inside one passage}
\label{sec:timing}

\begin{table}[t]
\centering
\caption{Timing budget for one passage at $15\,\mathrm{Hz}$ with five kept frames. The segmenter line is a budget, not a measured latency of a named device. The other lines are consequences of the algorithms in Chapter~\ref{ch:methods}.}
\label{tab:timing}
\small
\begin{tabular}{@{}lcc@{}}
\toprule
Stage & When & Budget \\
\midrule
Exposure and readout & every frame & $\le 2\,\mathrm{ms}$ exposure \\
YOLO-vf-seg & every kept frame & one frame period \\
ICPH, all masks of the frame & every kept frame & $<1\,\mathrm{ms}$ \\
Sharp-frame selection & trigger close & negligible \\
Homography on the centers & once & negligible \\
PLC origin, snap, completion & once & negligible \\
Write the tally record & once & negligible \\
\bottomrule
\end{tabular}
\end{table}

\Cref{tab:timing} is a budget that makes the dependence explicit. The system is late if and only if the segmenter is late, or if the trigger opens so late that fewer than one sharp frame is buffered. Both of those failures are detectable. A late segmenter drops a frame and still returns a tally from the frames it finished. An empty buffer returns no tally and an explicit failure, rather than a zero count. A zero count is a legal tally of an empty pallet and must not be reused as the code for ``the camera blinked.''

\subsection{The record the operator sees}
\label{sec:sys-output}

The record of a passage contains one line per color and one total, plus the diagnostics that change the meaning of those numbers. The fields are:

\begin{enumerate}
\item time, gate identity, and the frame index that PLC used;
\item $\hat{N}_{k}$ for each color, computed by~\eqref{eq:per-color}, with uncertain masks listed separately and not folded into a color;
\item the total $\hat{N}=\hat{d}\,\lvert\mathcal{F}\rvert$;
\item the depth reading and whether it was a default;
\item the number of completed cells $\lvert\mathcal{F}\setminus\mathcal{D}\rvert$;
\item the median origin residual, in centimetres;
\item the minimum color margin on the chosen frame.
\end{enumerate}

A total without the completed-cell count would hide the difference between a stack that was seen and a stack that was inferred. A per-color tally without the minimum margin would hide an amber--yellow swap. The seven fields are the smallest set that matches the failure modes of \Cref{fig:failures}. The display does not show a live mask overlay as a requirement. An overlay is useful during calibration and is not part of the tally the warehouse books.

Fault handling is a list of refusals. The system refuses to book a total when the buffer is empty, when the origin residual exceeds half a cell, or when $\hat{d}$ is missing and no shift default exists. In each refusal the images are kept and the passage is marked for a person. The system does not refuse merely because a margin is low. It books the total, withholds the uncertain color, and says so. That split follows from the independence of the two modules: a bad color is not a bad count.

\subsection{What the construction deliberately leaves out}
\label{sec:leave-out}

No second camera is specified. No truck-mounted computer is specified. No automatic white balance is permitted. No density branch is added to the segmenter. No multi-frame union is used in the tally. Each omission corresponds to a result or a requirement already stated. Listing them here is part of the construction, because a later implementer who adds a union ``for robustness'' would be reversing Section~\ref{sec:plc-passage} without a new measurement. The place to revisit any of these omissions is a field trial that measures the corresponding error, not a default setting in the process.

The construction is complete enough to implement and incomplete in the way any unfilmed design is incomplete. It does not name a lens, a cable, or a model of PC. Those are procurement choices inside the geometric and timing constraints above. The constraints are the design.

\section{Estimation}
\label{sec:estimate}

\subsection{Status of the numbers}
\label{sec:status}

Every number in this section comes from one instrument, \texttt{simulate\_estimation.py}, run at a fixed seed, writing \texttt{estimation\_results.json}. The instrument implements the observation models of Sections~\ref{sec:color-model} and~\ref{sec:count-model}. It does not read an image. It does not fine-tune a detector. A table below is therefore an estimate of a module under a stated generative model. Where the model is harsh, the estimate is harsh. Where the model omits a nuisance, the estimate is silent on that nuisance. The discussion in Section~\ref{sec:reading} says which silences matter.

The instrument has two independent blocks. The color block trains the linear ICPH of~\eqref{eq:linear} and tests it against raw cosine matching, gray-world, and shades-of-gray. The counting block does not look at color. It drops centers with a position-dependent probability, snaps the survivors, and scores six counters. Independence is a modeling choice already used in~\eqref{eq:bound}. It means the tables must not be multiplied together into a single invented accuracy.

\subsection{Observation model for color}
\label{sec:color-model}

A test patch is drawn by the same steps as training, with the severity held fixed so that a curve against illuminant strength exists. The severity values are $0.05$, $0.15$, $0.30$, $0.45$, and $0.60$. At each value, $2500$ patches are drawn, the class uniform on the six prototypes, the gain as in~\eqref{eq:gain} with that fixed $s$, the shadow, the shade ramp, the neighbor strip with probability $0.25$, and the pixel noise. Training used $6000$ patches with $s$ random in $[0.05,0.55]$ and $500$ full-batch steps at learning rate $0.35$ and $\lambda=10^{-3}$. Test patches are not in the training set; they are fresh draws, which for this generator is the relevant notion of ``unseen.''

The four decision rules are exactly those of Section~\ref{sec:icph-score}. Raw matching uses~\eqref{eq:raw} on the uncorrected mean. Gray-world corrects with $p=1$ and then matches. Shades-of-gray corrects with $p=6$ and then matches. ICPH computes~\eqref{eq:feature}, standardizes with the training moments, and applies the trained $W$ and $b$. Accuracy is the fraction of patches whose predicted class equals the class that generated the reflectance, before the illuminant was applied. A neighbor strip does not change the label. The label is the class of the face that owns the patch, not the class of the contaminating strip. That choice makes boundary error count against the classifier, which is the deployment meaning of a messy mask.

\subsection{Observation model for counting}
\label{sec:count-model}

A trial draws a stack and then a noisy view of it. The number of front rows is uniform on $\{2,3,4,5,6\}$, the number of columns uniform on $\{2,3,4\}$, and the depth uniform on $\{1,2,3\}$. Each front cell is occupied with probability $0.93$, and the stack is forced to contain at least one crate. A blur variable $\beta$ is uniform on $[0,1]$. The probability that an occupied cell is detected is
\begin{equation}
\pi(r)
= \operatorname{clip}\!\big(\pi_0(r)-0.18\,\beta,\, 0.35,\, 0.99\big),
\label{eq:detprob}
\end{equation}
with $\pi_0=0.70$ on the bottom row, $0.84$ on the top row, $0.96$ on a middle row, and $0.90$ if the stack has only one row. A detected center is the cell center plus independent Gaussian noise of standard deviation $0.035\,\mathrm{m}$ in $x$ and $0.028\,\mathrm{m}$ in $y$. A Poisson number of false centers, mean $0.35$, is thrown uniformly into a rectangle slightly larger than the stack. Depth is observed by~\eqref{eq:depth}. Five frames are drawn independently for each trial. Four thousand trials are run.

The six counters are the mask tally~\eqref{eq:mask-count}, the mask tally times depth, the area ratio~\eqref{eq:area}, the span~\eqref{eq:span}, PLC on the first frame, and PLC on the union of the five frames with vote threshold one. The area-ratio fraction and the span percentile match Section~\ref{sec:plc-baselines}. PLC matches \Cref{alg:plc}. The truth is~\eqref{eq:true-count}.

\subsection{Metrics}
\label{sec:est-metrics}

Color uses top-1 accuracy at each severity, plus, at the harshest severity, the per-class recall and the full confusion matrix of ICPH. Counting uses, for each counter, the mean absolute error, the root mean square error, the signed bias, the fraction of trials with an exact total, the mean absolute percentage error, and the ninetieth percentile of the absolute error. Errors are $\hat{N}-N$, so a negative bias is a counter that is short. The same absolute and squared errors are broken out by true depth, because a counter that ignores depth will look acceptable at depth one and unacceptable at depth three. The mean true count in the four thousand trials is $22.62$ with standard deviation $14.57$, which is the denominator that makes a mean absolute error of three crates a relative error of roughly one seventh, not a small or a large number by itself.

\subsection{Color results}
\label{sec:color-results}

\begin{table}[t]
\centering
\caption{Top-1 color accuracy on $2500$ patches at each illuminant severity. Gray-world uses $p=1$. Shades-of-gray uses $p=6$. ICPH is the linear head on the eight-dimensional feature. The generator, not a photographed aisle, produced the patches.}
\label{tab:color}
\small
\begin{tabular}{@{}lcccc@{}}
\toprule
Severity $s$ & Raw RGB & Gray-world & Shades-of-gray & ICPH \\
\midrule
$0.05$ & $0.986$ & $0.165$ & $0.196$ & $0.972$ \\
$0.15$ & $0.900$ & $0.180$ & $0.211$ & $0.914$ \\
$0.30$ & $0.778$ & $0.186$ & $0.219$ & $0.834$ \\
$0.45$ & $0.618$ & $0.172$ & $0.208$ & $0.687$ \\
$0.60$ & $0.531$ & $0.171$ & $0.204$ & $0.607$ \\
\bottomrule
\end{tabular}
\end{table}

\begin{figure}[t]
\centering
\begin{tikzpicture}
\begin{axis}[
  width=0.92\textwidth,
  height=7.4cm,
  xlabel={Illuminant severity $s$},
  ylabel={Top-1 accuracy},
  ymin=0, ymax=1.05,
  xmin=0, xmax=0.65,
  legend style={font=\small, at={(0.5,-0.22)}, anchor=north, legend columns=2},
  tick label style={font=\small},
  label style={font=\small},
  grid=major,
  grid style={gray!25}
]
\addplot+[mark=*] coordinates {(0.05,0.9864) (0.15,0.9000) (0.30,0.7780) (0.45,0.6184) (0.60,0.5312)};
\addplot+[mark=square*] coordinates {(0.05,0.1652) (0.15,0.1796) (0.30,0.1864) (0.45,0.1720) (0.60,0.1708)};
\addplot+[mark=triangle*] coordinates {(0.05,0.1956) (0.15,0.2112) (0.30,0.2188) (0.45,0.2080) (0.60,0.2040)};
\addplot+[mark=diamond*] coordinates {(0.05,0.9716) (0.15,0.9140) (0.30,0.8336) (0.45,0.6872) (0.60,0.6072)};
\legend{Raw RGB, Gray-world, Shades-of-gray, ICPH}
\end{axis}
\end{tikzpicture}
\caption{The same numbers as \Cref{tab:color}. Classical correction sits near chance for every severity. Raw matching wins only while the lamp is nearly white. ICPH degrades more slowly once the severity leaves that regime.}
\label{fig:color-curve}
\end{figure}

\Cref{tab:color} and \Cref{fig:color-curve} match the contract of Section~\ref{sec:icph-cost}. Gray-world and shades-of-gray stay between $0.165$ and $0.219$ at every severity. Chance on six classes is $0.167$. \Cref{prop:uniform} said a uniform face would be driven to neutral by the correction, and a face that is only mildly contaminated is not driven all the way, which is why the two classical rules sit just above chance rather than on it. They do not become more accurate as the illuminant grows, because the destruction of reflectance happens before the illuminant is interesting.

Raw matching at $s=0.05$ reaches $0.986$, above ICPH at $0.972$. That crossover was predicted: when the gain is almost white, dividing by $\hat{\ell}$ spends a clean reflectance signal to buy an illuminant correction that was unnecessary. By $s=0.15$ the order has reversed, $0.914$ against $0.900$. The gap then widens and stays modest: $0.056$ at $s=0.30$, $0.069$ at $s=0.45$, and $0.076$ at $s=0.60$. ICPH at the harshest setting is $0.607$, which is better than raw and far from solved. A deployment that sees lamps of this severity will abstain often, which is why the margin rule exists. The curve does not support a claim that the chromatic feature has removed illumination.

\begin{table}[t]
\centering
\caption{Confusion matrix of ICPH at severity $0.60$. Rows are the true class, columns the predicted class, in the order amber, green, red, blue, yellow, white. Each row sums to the number of test patches of that class inside the $2500$.}
\label{tab:confusion}
\small
\begin{tabular}{@{}lrrrrrr@{}}
\toprule
True $\backslash$ predicted & amber & green & red & blue & yellow & white \\
\midrule
amber  & 171 & 15 & 43 & 0 & 118 & 19 \\
green  & 4 & 271 & 2 & 51 & 33 & 46 \\
red    & 53 & 1 & 321 & 2 & 6 & 29 \\
blue   & 0 & 48 & 0 & 361 & 0 & 38 \\
yellow & 130 & 54 & 41 & 0 & 179 & 37 \\
white  & 21 & 66 & 34 & 70 & 21 & 215 \\
\bottomrule
\end{tabular}
\end{table}

\Cref{tab:confusion} is the harsh-severity detail behind the mean of $0.607$. Per-class recall is $0.467$ for amber, $0.666$ for green, $0.779$ for red, $0.808$ for blue, $0.406$ for yellow, and $0.504$ for white. Amber sends $118$ patches to yellow, and yellow sends $130$ patches to amber. That pair is the failure named in Section~\ref{sec:icph-cost}, and it is large enough that a mean accuracy which hides it would be a misleading summary. Red and blue are the classes the feature can hold. White's errors are spread across all five chromatic classes, $21$, $66$, $34$, $70$, and $21$, which is the neutral-axis prediction and not a single swap. Green leaks toward blue and white more than toward red, consistent with a prototype that already sits between those neighbors in \Cref{fig:chroma-plane}.

The operating consequence is narrower than ``use ICPH.'' Under a mild lamp, raw matching is enough and slightly better. Under a lamp the weekly calibration shows to be strong, ICPH is the better of the two and still needs the abstention rule, aimed first at amber and yellow. Under either lamp, a corrected-then-matched baseline is not a competitor. It is a demonstration of \Cref{prop:uniform}.

\subsection{Counting results}
\label{sec:count-results}

\begin{table}[t]
\centering
\caption{Counting error on $4000$ stacks. Bias is the mean of $\hat{N}-N$. Exact is the fraction of trials with $\hat{N}=N$. PLC-single uses the first frame. PLC-passage unions five frames at vote one. The mean true count is $22.62$.}
\label{tab:count}
\small
\begin{tabular}{@{}lrrrrrr@{}}
\toprule
Counter & MAE & RMSE & Bias & Exact & MAPE & P90 $\lvert e\rvert$ \\
\midrule
Mask, no depth & $13.55$ & $18.04$ & $-13.53$ & $0.044$ & $0.517$ & $31$ \\
Mask $\times$ depth & $4.87$ & $6.69$ & $-4.33$ & $0.126$ & $0.229$ & $12$ \\
Area ratio & $7.26$ & $9.19$ & $-7.08$ & $0.030$ & $0.328$ & $15$ \\
Span lattice & $3.31$ & $5.50$ & $+1.60$ & $0.332$ & $0.178$ & $9$ \\
PLC-single & $3.01$ & $5.00$ & $+0.83$ & $0.337$ & $0.162$ & $8$ \\
PLC-passage & $5.80$ & $8.79$ & $+5.33$ & $0.221$ & $0.318$ & $15$ \\
\bottomrule
\end{tabular}
\end{table}

\begin{figure}[t]
\centering
\begin{tikzpicture}
\begin{axis}[
  width=0.95\textwidth,
  height=7.2cm,
  ybar,
  bar width=7pt,
  ylabel={Mean absolute error (crates)},
  symbolic x coords={Mask,Mask$\times d$,Area,Span,PLC-1,PLC-F},
  xtick=data,
  x tick label style={font=\small},
  y tick label style={font=\small},
  ymin=0, ymax=16,
  nodes near coords,
  every node near coord/.append style={font=\scriptsize, /pgf/number format/fixed},
  label style={font=\small}
]
\addplot coordinates {(Mask,13.55) (Mask$\times d$,4.87) (Area,7.26) (Span,3.31) (PLC-1,3.01) (PLC-F,5.80)};
\end{axis}
\end{tikzpicture}
\caption{Mean absolute error of the six counters, from \Cref{tab:count}. The vertical axis is crates per passage, on a mean load of about $23$ crates. PLC on one frame is the lowest bar. The five-frame union is not.}
\label{fig:count-bars}
\end{figure}

\Cref{tab:count} is the counting claim of Section~\ref{sec:plc-claims}, measured. The mask tally without depth has bias $-13.53$ and mean absolute error $13.55$. The two numbers are almost equal, which means the error is nearly a pure under-count: the tally sees the front and the truth includes the ranks behind it. Multiplying the same masks by $\hat{d}$ cuts the mean absolute error to $4.87$ and the bias to $-4.33$. Depth is the dominant correction. Any discussion of lattice refinement that omits this comparison will attribute to geometry what belongs to the integer $\hat{d}$.

Given depth, the area ratio is worse than the raw mask count, $7.26$ against $4.87$, with bias $-7.08$. The visible-fraction model has mean $0.86$. Summing fractions that are systematically below one, then rounding, short-changes a stack whose detections were already a little short. The area ratio would need a fraction model centered on one, or a completion rule of its own, before it could be the inventory method of Hu, Wang, and Yang~\cite{hu2025arbca} inside this generator. The comparison does not refute that paper's measurement on tobacco cargo. It shows that a fraction below one is the wrong summary of a missed crate of known size.

The span and PLC-single are the two counters that use the pitch. Span reaches mean absolute error $3.31$, bias $+1.60$, exact on $0.332$ of trials. PLC-single reaches $3.01$, bias $+0.83$, exact on $0.337$. The absolute-error gap is three tenths of a crate on a load of about twenty-three. The bias gap is the clearer signature of the gap rule: both methods over-fill, and the two-neighbor rule over-fills less than the solid rectangle of the span. Neither is exact on a majority of trials. A ninetieth percentile of eight crates for PLC-single is a large tail. The tail is where the depth reading is off by one, because a wrong rank costs $\lvert\mathcal{F}\rvert$ crates at once, or where the outline is a step that completion fills.

PLC-passage is the negative result. Mean absolute error $5.80$, bias $+5.33$, exact on $0.221$. The bias and the absolute error are close, so the passage union is a nearly pure over-count, as Section~\ref{sec:plc-passage} argued from the vote threshold of one. Five frames are worse than one frame. The system specification that runs PLC once is this row of the table, not a preference for simpler code.

\begin{table}[t]
\centering
\caption{Mean absolute error by true depth. The number of trials is $1326$, $1329$, and $1345$ at depths $1$, $2$, and $3$.}
\label{tab:bydepth}
\small
\begin{tabular}{@{}lrrr@{}}
\toprule
Counter & Depth $1$ & Depth $2$ & Depth $3$ \\
\midrule
Mask, no depth & $2.27$ & $13.30$ & $24.93$ \\
Mask $\times$ depth & $2.50$ & $4.89$ & $7.18$ \\
Area ratio & $3.53$ & $7.22$ & $10.98$ \\
Span lattice & $1.91$ & $3.59$ & $4.40$ \\
PLC-single & $1.72$ & $3.27$ & $4.03$ \\
PLC-passage & $3.29$ & $5.96$ & $8.12$ \\
\bottomrule
\end{tabular}
\end{table}

\Cref{tab:bydepth} separates the depth effect. The mask tally's error scales almost exactly with the hidden ranks: $2.27$, $13.30$, $24.93$. At depth one it is competitive with the geometric counters, because there is no hidden rank and the remaining error is detection noise. At depth three it is an inventory error of a different kind. Mask times depth grows much more slowly, $2.50$, $4.89$, $7.18$, and the growth is the combination of front-face misses scaled by $\hat{d}$ and the ten percent of trials in which $\hat{d}$ itself is wrong. PLC-single is the best entry in every column: $1.72$, $3.27$, $4.03$. The margin over the span is small and has the same sign in every column, which is what a refinement rather than a different estimator looks like. PLC-passage is worse than PLC-single in every column, so the over-fill is not an artifact of one depth.

Bias, not shown in that table and computed in the same run, stays negative for the mask and the area ratio at every depth, and stays positive for the span and both PLC variants at every depth. The sign flip between ``count what you saw'' and ``complete a lattice'' is stable across the design range. An operator who trusts PLC is accepting a small positive bias in exchange for a smaller absolute error. The completed-cell field on the record is the way that acceptance is made visible per passage rather than only in a table of four thousand trials.

\subsection{One passage, computed by hand and by the instrument}
\label{sec:one-passage}

One draw in the run has the front occupancy~\eqref{eq:example-grid}, depth $2$, and therefore truth $20$. The blur draw is $0.161$. The five frames contain $9$, $10$, $10$, $10$, and $10$ detections. The first frame is one short of the ten true front cells. PLC on that frame returns $22$. The span on that frame returns $24$. The mask tally returns $9$, which is the front count with one miss and no depth. PLC-passage on all five frames returns $28$.

Section~\ref{sec:plc-example} already accounts for the $22$: one missed cell put back by the neighbor rule, one phantom cell added on the stepped top row, eleven front cells times depth $2$. The span's $24$ is the bounding rectangle, three columns by four rows, times depth $2$, which fills both empty cells of the top row rather than one of them. The passage result of $28$ is fourteen front cells times depth $2$, four cells above the truth on the front face. Those four are the union's contribution. The instrument and the hand rule agree on the frame that was short, which is the check that \Cref{alg:plc} in the text and the instrument are the same algorithm. They do not agree with the truth. The example is a picture of the positive bias, chosen because the grid is small enough to count, not because the error is flattering.

\subsection{Parameter cost of the head}
\label{sec:cost-result}

The cost model of~\eqref{eq:params} is a count of weights, not a profiled runtime. The reference segmenter in that model is a YOLOv8m-seg scale network, $27.3$ million parameters and on the order of $10^{2}$ GFLOPs, whose published COCO box and mask figures sit near $50$ and $41$ mean average precision. Those COCO figures are a reference for the size of the network. They are not a prediction of gate accuracy. On top of that network the specified ICPH adds $9760$ parameters, a fraction of $3.6\times 10^{-4}$. The counting stage adds none. The estimate of accuracy and the estimate of cost answer different questions. The cost answer is that the color head is not the part of the system a deployer would shrink first. The accuracy answer is that the head is worth having when the illuminant is not mild, and that it does not resolve amber against yellow at severity $0.60$.

\subsection{How to read the estimate}
\label{sec:reading}

The estimate supports five statements, and it does not support a sixth.

It supports the statement that classical per-mask color constancy is the wrong correction on a nearly uniform crate face. That statement is \Cref{prop:uniform} plus the flat curves near chance in \Cref{fig:color-curve}. It does not depend on the linear head being optimal.

It supports the statement that a linear head on $(\chi,o,\eta)$ degrades more slowly than raw cosine matching as the illuminant severity grows, and that the residual errors concentrate on amber, yellow, and white. The crossover at low severity is part of the statement, not a footnote.

It supports the statement that multiplying a front count by a depth reading is the largest single reduction in counting error available inside this model. The reduction from $13.55$ to $4.87$ mean absolute error dominates the later geometric refinements.

It supports the statement that, given depth, a lattice snap with two-neighbor completion is a small improvement on a span fill and a clear improvement on an area ratio whose visible fraction averages below one. The improvement is a few tenths of a crate on the mean, and a reduction of positive bias.

It supports the statement that unioning five frames at a vote of one increases positive bias. The recommended system runs PLC on one frame because of this statement, not in spite of it.

It does not support a field accuracy, a claim about cardboard of unknown pitch, a claim about a different depth-sensor error rate, or a claim that the neck path of Section~\ref{sec:icph-neck} has been trained. Those would require images, a second pitch, a measured sensor, and a trained backbone. The instrument does not contain them, and the tables do not imply them.

A reader who needs a single comparison should take it from \Cref{tab:count} and \Cref{tab:color} together and keep the conditions attached. At severity $0.30$, ICPH names the design palette correctly on about five patches in six, against about four in five for raw matching and about one in five for a classical correction. On the counting model, PLC-single is off by about three crates on a load of about twenty-three, against about five for masks times depth and about fourteen for masks alone. Both sentences are about the generators in Sections~\ref{sec:color-model} and~\ref{sec:count-model}.

\subsection{What a field trial would have to add}
\label{sec:field}

A trial that could replace this estimate would record, at one gate, the images, the depth reading, a person-confirmed count per color, and the lamp condition in enough variety that severity has an empirical meaning. It would report the same metrics: signed bias, exact total, exact per-color count, and abstention rate. It would keep PLC-single and PLC-passage as separate rows, because the passage result is easy to ``improve'' by a code change that the model has already rejected. It would publish the confusion matrix, not only the mean, because amber and yellow can move the mean by a point while dominating the inventory complaints.

Until that trial exists, the system of Section~\ref{sec:build} is a specification and the tables of this section are an estimate of its two numerical modules. The specification is already constrained by the estimate: one frame, not five; a chromatic feature, not a corrected color; depth as an input, not as a fraction of mask area. That is the sense in which the estimation chapter closes the design. It does not close the measurement.

\subsection{Footprint of a pixel on the crate face}
\label{sec:footprint}

The timing budget and the snap radius should be checked against the size of a pixel, otherwise the centimetre noise in the observation model is an unattached number. Take the installation of Section~\ref{sec:gate}: standoff $Z=2.5\,\mathrm{m}$, face width of interest $1.6\,\mathrm{m}$ for four crates, and a sensor used in landscape orientation. A horizontal field of view of $50^{\circ}$ covers
\begin{equation}
2\,Z\tan(25^{\circ}) = 2\cdot 2.5\cdot 0.466 = 2.33\,\mathrm{m}
\label{eq:fov}
\end{equation}
at the front plane, which leaves margin around a $1.6\,\mathrm{m}$ face for a truck that is not perfectly centered. Across $1280$ horizontal pixels the footprint is
\begin{equation}
\frac{2.33}{1280} = 1.8\,\mathrm{mm}
\label{eq:footprint}
\end{equation}
per pixel at the center of the field. The vertical footprint is comparable near the camera height and larger toward the top and bottom rows, where the ray is oblique. A center error of one pixel is negligible beside the $3.5\,\mathrm{cm}$ localization noise of the model and beside the $20\,\mathrm{cm}$ half-cell of \Cref{prop:snap}. The lattice is not limited by sensor resolution at this standoff. It is limited by which crates are detected and by whether the homography scale is right.

The same footprint converts the exposure specification into pixels of smear. At $1.5\,\mathrm{m\,s^{-1}}$ and $2\,\mathrm{ms}$ the load moves $3\,\mathrm{mm}$, about two pixels. At $5\,\mathrm{m\,s^{-1}}$ and $2\,\mathrm{ms}$ it moves $10\,\mathrm{mm}$, about six pixels. Six pixels of horizontal smear on a vertical lip, which is the lip the vertical-feature gate is trying to see, is a loss of contrast rather than a displacement of the center by a full cell. The $2\,\mathrm{ms}$ exposure is the specification that keeps smear in that regime. A camera that cannot reach it at the illumination of the dock has to add light. The design does not respond to a dark dock by lengthening the shutter, because the snap radius is a counting requirement and the brightness is a procurement requirement.

\begin{figure}[t]
\centering
\begin{tikzpicture}[x=0.72cm,y=0.72cm,font=\small]
  \draw[thick] (0,0) -- (14,0);
  \draw[dashed] (2,0) -- (2,6.2) node[above] {\footnotesize truck path};
  \draw[dashed] (8.2,0) -- (8.2,6.2);
  \node[align=center] at (5.1,5.6) {\footnotesize front-face plane};
  \draw[{Stealth}-{Stealth}] (2,5.1) -- (8.2,5.1) node[midway, above] {$2.5\,\mathrm{m}$};
  \draw[fill=black!80] (10.6,2.4) rectangle (11.7,3.1);
  \draw[-{Stealth}, thick] (10.6,2.75) -- (8.2,2.75);
  \node[align=center] at (11.15,3.55) {\footnotesize camera\\[-1pt]\footnotesize $1.4\,\mathrm{m}$};
  \draw[{Stealth}-{Stealth}] (12.3,0) -- (12.3,2.75) node[midway, right] {\footnotesize $1.4\,\mathrm{m}$};
  \draw[fill=orange!30] (4.3,0.4) rectangle (5.9,3.6);
  \node[align=center] at (5.1,2.0) {\footnotesize stack};
  \node[align=center] at (5.1,-0.55) {\footnotesize gate floor};
\end{tikzpicture}
\caption{Plan-and-elevation sketch of the gate, not to a common scale. The camera is fixed, the face plane is the calibration plane of $H$, and the standoff is the $Z$ used in~\eqref{eq:fov}. The truck path is perpendicular to the optical axis at the moment PLC is meant to run.}
\label{fig:gate}
\end{figure}

\subsection{From pixels to a booked line}
\label{sec:booked-line}

The pieces are scattered through the derivations. One passage, with the numbers of the worked example, shows the order in which the computer applies them. The trigger opens. Five frames are buffered. On each frame YOLO-vf-seg emits masks. ICPH scores each mask. Suppose the first frame has the lowest blur and the highest median margin, and that its ten true front crates produced nine detections, as in the instrument's draw. The nine centers are multiplied by $H$ and snapped. The origin residual is a few centimetres, inside the half-cell, so the passage is not refused. The gap rule adds the restored miss and the one phantom cell of \Cref{fig:grid}. The depth reading returns $2$, which matches the truth in this particular draw. The total booked is $11\times 2=22$ against a truth of $20$. The completed-cell count is $2$ if both the restored cell and the phantom were absent from the raw snap, and the record shows that $2$. The color head names the nine detected faces. The two cells that were not detected inherit by~\eqref{eq:per-color}. If one of the nine was an amber face scored as yellow with a margin below the weekly tenth percentile, that face is listed as uncertain and is not forced into either color's tally, while the total stays $22$.

The operator who sees this record can do three different things, and the record is built so that they stay different. A total of $22$ with two completed cells and a small residual is a geometric inference, to be trusted to the extent the building trusts PLC-single's bias of under one crate on average. An uncertain amber is a color question, to be confirmed from the image or from the pallet ticket, without reopening the total. A depth field that says ``default'' rather than ``from the pallet record'' is a rank question. Collapsing all three into a red flag would send every stepped stack to a person. Separating them sends the person the question the model cannot answer.

The same passage under the counters that were rejected reads as follows. The mask tally books $9$ and is short by $11$, because it ignores depth and one miss. Mask times depth books $18$ and is short by $2$. The span books $24$, the full rectangle times depth, and is long by $4$. PLC-passage, given the instrument's five detection counts of $9,10,10,10,10$, books $28$ and is long by $8$. The booked line of the specified system is the least wrong of those integers on this draw, and it is still wrong by two crates. The worked passage is not a demonstration that the error has been removed. It is a demonstration that the record contains the integer and the reason the integer might be high.

\subsection{Sensitivity to the depth reading}
\label{sec:depth-sensitivity}

The depth model fixes a ten percent chance of being off by one. The tables do not vary that chance, and a reader will want the direction of the dependence. Condition on the front estimate $\lvert\mathcal{F}\rvert$ being some integer $m$, and on the truth $N=d\,m_{\star}$. When $\hat{d}=d$, the counting error is $d(m-m_{\star})$, the front error scaled by depth. When $\hat{d}=d+1$, the error gains an extra $m$. When $\hat{d}=d-1$, it loses $m$. At the stated probability $0.10$ of a wrong reading, split equally between the two directions and ignoring the clip at the ends of $\{1,2,3\}$, the expected contribution of depth noise alone is of order $0.1\,m$ in absolute error and near zero in bias, because the two directions cancel in the signed mean and add in the absolute mean. On a front of about ten cells, depth noise contributes on the order of one crate of mean absolute error. That is a substantial part of PLC-single's $3.01$, and it is a part no lattice rule can remove.

Two consequences follow for the building, both outside the Monte Carlo grid and both readable from the scaling. Halving the depth error rate, from $0.10$ to $0.05$, is worth on the order of half a crate of mean absolute error if the front estimate stays as it is. That is comparable to the entire gap between the span and PLC-single. A better depth source can therefore matter as much as the geometric refinement this thesis spent a section on. Improving the front estimate without improving $\hat{d}$ still matters, because front errors are multiplied by $d$, which is why PLC-single's advantage over the span grows from depth $1$ to depth $3$ in \Cref{tab:bydepth} even though the rule itself does not see $d$. The design priority that follows is: obtain $\hat{d}$ from the pallet record whenever the record exists, and spend the vision effort on the front face. Spending the vision effort on inventing $\hat{d}$ from the side image would put a high-variance estimate into the factor that multiplies every other error.

\subsection{What the estimate constrains in the build}
\label{sec:constraints}

The specification in Section~\ref{sec:build} was written to be consistent with the tables, and the consistencies are checkable.

The camera's white balance is locked because gray-world and shades-of-gray, including any global gain that behaves like them on a uniform face, destroy the type. The lock is not a preference for manual color. It is \Cref{fig:color-curve}.

PLC runs on one frame because the five-frame union at vote one has bias $+5.33$. A software default that fused frames would contradict \Cref{tab:count}.

The total is multiplied by a depth reading because the mask tally without that multiplication has bias $-13.53$. A vision-only total is the front count.

The record shows completed cells and the minimum margin because PLC-single's remaining error is a small positive bias from cells the rule inferred, and ICPH's remaining error is a swap of amber and yellow rather than a uniform sprinkling of mistakes. A record that showed only $\hat{N}$ would hide both.

The color head is not the place the implementation looks for speed, because $9760$ parameters do not move the budget in \Cref{tab:timing}. The segmenter is that place.

None of these constraints is a field result. Each one is a decision that would be wasteful to relitigate inside the same observation model. A field trial can overturn any of them. Until it does, they are the decisions the estimate is for.

\subsection{Trace from requirement to row of the table}
\label{sec:trace}

The requirements of Section~\ref{sec:requirements} were a filter on the design. The estimate is now in hand, and the filter can be read in the other direction: each requirement is either discharged by a row of a table, or it is a property of the specification that no table can discharge. \Cref{tab:trace} is that reading. It is the checklist a later implementer can use without re-deriving the thesis. A requirement marked ``specification'' is discharged by the absence of a code path, not by a number. A requirement marked by a table is discharged only inside the observation model of that table.

\begin{table}[t]
\centering
\caption{Where each requirement of Section~\ref{sec:requirements} is answered.}
\label{tab:trace}
\small
\begin{tabular}{@{}p{0.28\textwidth}p{0.30\textwidth}p{0.32\textwidth}@{}}
\toprule
Requirement & Where it is answered & What would falsify it \\
\midrule
Type and total & \Cref{eq:per-color}, \Cref{tab:color}, \Cref{tab:count} & A joint trial whose per-color error exceeds~\eqref{eq:bound} in a systematic way \\
One passage & PLC-single row of \Cref{tab:count}; the passage row is the rejected alternative & A vote rule whose bias beats PLC-single on filmed passages \\
Side view only & Depth is an input in \Cref{alg:plc} and in Section~\ref{sec:compute} & A side-image depth estimator with smaller error than the pallet record \\
Illumination unlocked & \Cref{tab:color} and \Cref{tab:ablation}; firmware white balance forbidden & A lamp distribution under which raw matching dominates at every severity the dock produces \\
The seam & Vertical gate~\eqref{eq:vf-again}; not estimated & Labeled frames where merged seams dominate the residual of PLC \\
Line budget & \Cref{tab:timing}, $9760$ parameters & A profiled segmenter that misses the frame period after the head is added \\
Signed errors & Bias column of \Cref{tab:count} & An operation in which positive and negative inventory errors have the same cost and the same cause \\
\bottomrule
\end{tabular}
\end{table}

Two rows of \Cref{tab:trace} are the ones most likely to be overturned by a real dock, and they should be named rather than buried. The seam is not estimated. The Monte Carlo center noise is an average of many ways a center can move, including a merged seam that slides the moment~\eqref{eq:moment} halfway between two crates. If real frames merge seams often enough that the center falls outside the snap radius, PLC will book one cell where two crates stand, the completed-cell count will not rise, and the bias formula~\eqref{eq:front-signed} will acquire a new negative term the present $q$ does not represent, because the miss will be invisible to a counter that never received two centers. The vertical gate is the specified mitigation. It is untested. A field trial that labels seams is the trial that can retire or confirm that part of the architecture.

The illumination row can be overturned in the other direction. If the dock's lamps, after the camera gain is locked, stay inside the mild end of \Cref{tab:color}, raw matching is the better head and ICPH's correction is a small loss. The weekly margin calibration will show this as a high abstention rate that buys nothing, or as a confusion matrix that is already diagonal under the raw rule. The specified response is to keep the feature and the linear layer, because they are cheap, and to lower the influence of the correction only if a refit of $W$ on that building's patches puts near-zero weight on the residual coordinates. The ablation says the weights that matter are the illuminant block and the residual chromaticity. A refit is allowed to discover that, on this dock, the illuminant block alone is enough. The refit is not allowed to replace the head by a gray-world division. That alternative is not close.

\subsection{What a passage stores}
\label{sec:stores}

The record of Section~\ref{sec:sys-output} is what a person reads. The computer should store a little more, because the estimate cannot be replaced by a field trial that did not keep the inputs of the estimate. For each booked passage the store keeps the chosen frame, the masks of that frame, the centers in pixels and in metres, the origin residual, the set $\mathcal{D}$ before completion, the set $\mathcal{F}$ after completion, the depth reading and its source bit, the per-mask class and margin, and the seven displayed fields. It does not keep every frame of the buffer once the sharp frame has been chosen, except when the passage was refused, in which case the whole buffer is kept. The distinction is deliberate. A refused passage is the data from which a better trigger or a better shutter will be judged. A booked passage is the data from which PLC and ICPH will be judged, and the unused frames are the temptation to reintroduce the union that \Cref{tab:count} rejected.

Retention is a shift, not a year. The geometric calibration target images are kept until the next calibration. The weekly color-calibration patches are kept until the next refit of $W$, together with the vector $W$ itself, so that a change in amber--yellow recall can be attributed to the weights or to the lamps. None of this store is a contribution of the thesis. It is the minimum on which the claim ``the tables were replaced by a measurement'' can later be true. A deployment that books $\hat N$ and discards the centers cannot compute the origin residual after the fact, cannot tell a phantom cell from a detection, and cannot rebuild \Cref{tab:confusion} when a customer complains about yellow crates booked as amber.

The size of the store is small next to the video the camera could have written. One chosen frame at the working resolution, a few dozen masks as compressed bitmaps, and a page of numbers, times the number of trucks in a shift, is an ordinary disk. The constraint is not capacity. The constraint is that the booked integer and the evidence for the integer stay in the same record, so that the signed error the thesis treats as meaningful can be computed on a Friday from the passages of the week.

\subsection{What the mean error is, in a shift}
\label{sec:shift}

A mean absolute error of three crates is an abstract number until it is placed against a shift. The counting model drew loads whose mean truth was $22.6$ crates, so three crates is about thirteen percent of a typical load and, on the depth-three slice, four crates against a larger load. Suppose a shift sends forty pallets through the gate, independent in the sense of the model, and each is booked by PLC-single. The expected total absolute discrepancy across the shift is about $40\times 3.01 \approx 120$ crates, while the expected signed discrepancy is about $40\times 0.83 \approx 33$ crates long. The signed number is the one that survives in a daily reconciliation against the packaging line's own tally: misses and phantoms that cancel inside one pallet do not cancel in the sum of absolute errors, and the bias adds. A warehouse that ``checks the system'' by comparing the day's sum to the day's production will see a discrepancy near the bias, and will conclude that the system is long by a fraction of a crate per pallet. That conclusion is the bias column. It is not the mean absolute error, and it is not a reason to turn the gap rule off. Turning the rule off returns the system to the mask-times-depth row, whose bias is $-4.33$ per pallet, about $170$ crates short over the same forty pallets, a reconciliation error several times larger and of the dangerous sign.

The color errors do not add in the same way. A swap of amber and yellow preserves the total crates and moves one pallet's line from one stock-keeping unit to another. Over forty pallets, at the harsh severity where amber recall is $0.47$ and yellow recall is $0.41$, a large minority of those two colors would be booked under the other name if the head were forced to answer. The abstention rule exists so that this minority is a list of faces for a person rather than a silent transfer between two bins in the inventory. The shift arithmetic for color is therefore a count of abstentions, not a count of crates. The store of Section~\ref{sec:stores} keeps the margins so that the tenth-percentile threshold can be recomputed from the week's own pallets rather than from the generator's severity dial. Until that week exists, the dial and the shift arithmetic above are the scale on which the tables should be read: a few crates of absolute error per pallet, a bias under one crate and positive, and a color head that is asked to abstain on the pair whose prototypes sit $9.3^{\circ}$ apart.

The same shift makes the abstention rule quantitative. Suppose a pallet shows eight amber or yellow faces on the front, the harsh recall of those two classes is near one half, and the other colors are held well enough that the person is not asked about them. Four faces a shift, times forty pallets, is on the order of a hundred faces a day if every lamp is as harsh as $s=0.60$ and if every face is forced into the confusion rates of \Cref{tab:confusion}. That is not an operating point the design accepts. It is the upper end of the generator, and it is the reason the camera gain is locked and the reason the weekly calibration is allowed to move the abstention threshold. At the mild end of \Cref{tab:color} the same eight faces are named correctly almost all of the time by either ICPH or the raw matcher, and the list for the person is short. A dock that measures its own severity, by the margin distribution on the known-color pallets, will know which of those two days it is having. The thesis cannot know that for a building it has not seen. It can refuse to hide the range.

\begin{figure}[p]
\centering
\begin{tikzpicture}[
  font=\small,
  box/.style={draw, rounded corners=2pt, align=center, inner sep=4pt, minimum width=38mm, minimum height=11mm},
  arr/.style={-{Stealth}, thick}
]
  \node[box, fill=black!6] (cam) {Locked camera\\at the gate};
  \node[box, fill=orange!12, right=12mm of cam] (pass) {One passage\\five or fewer frames};
  \node[box, fill=orange!12, right=12mm of pass] (sharp) {Sharpest frame\\by median margin};
  \node[box, fill=green!12, below=12mm of sharp] (plc) {PLC-single\\and depth $\hat d$};
  \node[box, fill=blue!10, left=12mm of plc] (rec) {Record:\\total, colors,\\completed cells,\\weakest margin};
  \node[box, fill=black!5, left=12mm of rec] (week) {Weekly color\\calibration};
  \node[box, fill=red!8, below=12mm of rec] (abs) {Abstain on\\low-margin faces};
  \node[box, fill=black!6, below=12mm of plc] (shift) {Shift sum\\against the line};
  \draw[arr] (cam) -- (pass);
  \draw[arr] (pass) -- (sharp);
  \draw[arr] (sharp) -- (plc);
  \draw[arr] (plc) -- (rec);
  \draw[arr] (week) -- (rec);
  \draw[arr] (rec) -- (abs);
  \draw[arr] (rec) -- (shift);
\end{tikzpicture}
\caption{The passage the specification actually runs, collected onto one page so the estimate and the build can be read together. Frames are not unioned. Depth joins at PLC. The weekly calibration moves only the abstention threshold. The shift sum is where a positive bias of under one crate per pallet becomes a reconciliation against production, and where the much larger negative bias of a mask tally would appear if the lattice were removed.}
\label{fig:shift-flow}
\end{figure}


\subsection{Summary of the thesis}
\label{sec:summary}

The passage of a loaded forklift is a counting problem with a known pitch, a small integer depth, and a color inventory that has to survive the lamp. Published segmenters solve the visible mask and leave the color and the depth. Published constancy algorithms estimate a lamp for a mixed scene and erase a uniform face. Published area-ratio and volumetric counters recover hidden goods without using a crate lattice that the palletizer already imposed.

This thesis updates YOLO-vf-seg with two modules sized to those facts. ICPH keeps reflectance and illuminant as separate coordinates of one per-mask feature and scores six prototypes, at a cost of $9760$ parameters on the specified head. PLC snaps rectified centers to the crate pitch, fills a hole that has two orthogonal neighbors, and multiplies by a depth reading that the side image is not allowed to invent. The Monte Carlo estimate is consistent with the derivations. Classical correction falls to chance on a uniform face. ICPH holds a modest lead over raw matching once the illuminant leaves the white point, and it still confuses amber with yellow when the illuminant is harsh. Depth multiplication removes most of the counting error. The lattice removes a further, smaller share and will over-fill if it is allowed to union every frame of the passage.

The system that follows from those results is a gate camera with locked white balance, a segmenter that produces masks and centers, a color head that is allowed to abstain, and a lattice that runs once on the sharpest frame. The tally records the total, the colors, the depth reading, the number of cells that were inferred, and the weakest color margin. That record is the output the problem statement asked for. The numbers attached to it in this chapter are estimates under a model that the thesis writes down, and they remain estimates until a gate is filmed.

\clearpage


\bibliographystyle{unsrtnat}
\bibliography{refs}

\end{document}
